\documentclass[11pt]{article}

\usepackage[breakable]{tcolorbox}
\usepackage{parskip} % Stop auto-indenting (to mimic markdown behaviour)
\usepackage{ctex}
\usepackage[table]{xcolor}

% Basic figure setup, for now with no caption control since it's done
% automatically by Pandoc (which extracts ![](path) syntax from Markdown).
\usepackage{graphicx}
% Keep aspect ratio if custom image width or height is specified
\setkeys{Gin}{keepaspectratio}
% Maintain compatibility with old templates. Remove in nbconvert 6.0
\let\Oldincludegraphics\includegraphics
% Ensure that by default, figures have no caption (until we provide a
% proper Figure object with a Caption API and a way to capture that
% in the conversion process - todo).
\usepackage{caption}
\DeclareCaptionFormat{nocaption}{}
\captionsetup{format=nocaption,aboveskip=0pt,belowskip=0pt}

\usepackage{float}
\floatplacement{figure}{H} % forces figures to be placed at the correct location
\usepackage{xcolor} % Allow colors to be defined
\usepackage{enumerate} % Needed for markdown enumerations to work
\usepackage{geometry} % Used to adjust the document margins
\usepackage{amsmath} % Equations
\usepackage{amssymb} % Equations
\usepackage{textcomp} % defines textquotesingle
% Hack from http://tex.stackexchange.com/a/47451/13684:
\AtBeginDocument{%
	\def\PYZsq{\textquotesingle}% Upright quotes in Pygmentized code
}
\usepackage{upquote} % Upright quotes for verbatim code
\usepackage{eurosym} % defines \euro

\usepackage{iftex}
\ifPDFTeX
\usepackage[T1]{fontenc}
\IfFileExists{alphabeta.sty}{
	\usepackage{alphabeta}
}{
	\usepackage[mathletters]{ucs}
	\usepackage[utf8x]{inputenc}
}
\else
\usepackage{fontspec}
\usepackage{unicode-math}
\fi

\usepackage{fancyvrb} % verbatim replacement that allows latex
\usepackage{grffile} % extends the file name processing of package graphics
% to support a larger range
\makeatletter % fix for old versions of grffile with XeLaTeX
\@ifpackagelater{grffile}{2019/11/01}
{
	% Do nothing on new versions
}
{
	\def\Gread@@xetex#1{%
		\IfFileExists{"\Gin@base".bb}%
		{\Gread@eps{\Gin@base.bb}}%
		{\Gread@@xetex@aux#1}%
	}
}
\makeatother
\usepackage[Export]{adjustbox} % Used to constrain images to a maximum size
\adjustboxset{max size={0.9\linewidth}{0.9\paperheight}}

% The hyperref package gives us a pdf with properly built
% internal navigation ('pdf bookmarks' for the table of contents,
% internal cross-reference links, web links for URLs, etc.)
\usepackage{hyperref}
% The default LaTeX title has an obnoxious amount of whitespace. By default,
% titling removes some of it. It also provides customization options.
\usepackage{titling}
\usepackage{longtable} % longtable support required by pandoc >1.10
\usepackage{booktabs}  % table support for pandoc > 1.12.2
\usepackage{array}     % table support for pandoc >= 2.11.3
\usepackage{calc}      % table minipage width calculation for pandoc >= 2.11.1
\usepackage[inline]{enumitem} % IRkernel/repr support (it uses the enumerate* environment)
\usepackage[normalem]{ulem} % ulem is needed to support strikethroughs (\sout)
% normalem makes italics be italics, not underlines
\usepackage{soul}      % strikethrough (\st) support for pandoc >= 3.0.0
\usepackage{mathrsfs}


\usepackage{geometry}
\usepackage{float}
\usepackage{multicol}
\usepackage{subfigure}
\usepackage{makecell}

% Colors for the hyperref package
\definecolor{urlcolor}{rgb}{0,.145,.698}
\definecolor{linkcolor}{rgb}{.71,0.21,0.01}
\definecolor{citecolor}{rgb}{.12,.54,.11}

% ANSI colors
\definecolor{ansi-black}{HTML}{3E424D}
\definecolor{ansi-black-intense}{HTML}{282C36}
\definecolor{ansi-red}{HTML}{E75C58}
\definecolor{ansi-red-intense}{HTML}{B22B31}
\definecolor{ansi-green}{HTML}{00A250}
\definecolor{ansi-green-intense}{HTML}{007427}
\definecolor{ansi-yellow}{HTML}{DDB62B}
\definecolor{ansi-yellow-intense}{HTML}{B27D12}
\definecolor{ansi-blue}{HTML}{208FFB}
\definecolor{ansi-blue-intense}{HTML}{0065CA}
\definecolor{ansi-magenta}{HTML}{D160C4}
\definecolor{ansi-magenta-intense}{HTML}{A03196}
\definecolor{ansi-cyan}{HTML}{60C6C8}
\definecolor{ansi-cyan-intense}{HTML}{258F8F}
\definecolor{ansi-white}{HTML}{C5C1B4}
\definecolor{ansi-white-intense}{HTML}{A1A6B2}
\definecolor{ansi-default-inverse-fg}{HTML}{FFFFFF}
\definecolor{ansi-default-inverse-bg}{HTML}{000000}

% common color for the border for error outputs.
\definecolor{outerrorbackground}{HTML}{FFDFDF}

% commands and environments needed by pandoc snippets
% extracted from the output of `pandoc -s`
\providecommand{\tightlist}{%
	\setlength{\itemsep}{0pt}\setlength{\parskip}{0pt}}
\DefineVerbatimEnvironment{Highlighting}{Verbatim}{commandchars=\\\{\}}
% Add ',fontsize=\small' for more characters per line
\newenvironment{Shaded}{}{}
\newcommand{\KeywordTok}[1]{\textcolor[rgb]{0.00,0.44,0.13}{\textbf{{#1}}}}
\newcommand{\DataTypeTok}[1]{\textcolor[rgb]{0.56,0.13,0.00}{{#1}}}
\newcommand{\DecValTok}[1]{\textcolor[rgb]{0.25,0.63,0.44}{{#1}}}
\newcommand{\BaseNTok}[1]{\textcolor[rgb]{0.25,0.63,0.44}{{#1}}}
\newcommand{\FloatTok}[1]{\textcolor[rgb]{0.25,0.63,0.44}{{#1}}}
\newcommand{\CharTok}[1]{\textcolor[rgb]{0.25,0.44,0.63}{{#1}}}
\newcommand{\StringTok}[1]{\textcolor[rgb]{0.25,0.44,0.63}{{#1}}}
\newcommand{\CommentTok}[1]{\textcolor[rgb]{0.38,0.63,0.69}{\textit{{#1}}}}
\newcommand{\OtherTok}[1]{\textcolor[rgb]{0.00,0.44,0.13}{{#1}}}
\newcommand{\AlertTok}[1]{\textcolor[rgb]{1.00,0.00,0.00}{\textbf{{#1}}}}
\newcommand{\FunctionTok}[1]{\textcolor[rgb]{0.02,0.16,0.49}{{#1}}}
\newcommand{\RegionMarkerTok}[1]{{#1}}
\newcommand{\ErrorTok}[1]{\textcolor[rgb]{1.00,0.00,0.00}{\textbf{{#1}}}}
\newcommand{\NormalTok}[1]{{#1}}

% Additional commands for more recent versions of Pandoc
\newcommand{\ConstantTok}[1]{\textcolor[rgb]{0.53,0.00,0.00}{{#1}}}
\newcommand{\SpecialCharTok}[1]{\textcolor[rgb]{0.25,0.44,0.63}{{#1}}}
\newcommand{\VerbatimStringTok}[1]{\textcolor[rgb]{0.25,0.44,0.63}{{#1}}}
\newcommand{\SpecialStringTok}[1]{\textcolor[rgb]{0.73,0.40,0.53}{{#1}}}
\newcommand{\ImportTok}[1]{{#1}}
\newcommand{\DocumentationTok}[1]{\textcolor[rgb]{0.73,0.13,0.13}{\textit{{#1}}}}
\newcommand{\AnnotationTok}[1]{\textcolor[rgb]{0.38,0.63,0.69}{\textbf{\textit{{#1}}}}}
\newcommand{\CommentVarTok}[1]{\textcolor[rgb]{0.38,0.63,0.69}{\textbf{\textit{{#1}}}}}
\newcommand{\VariableTok}[1]{\textcolor[rgb]{0.10,0.09,0.49}{{#1}}}
\newcommand{\ControlFlowTok}[1]{\textcolor[rgb]{0.00,0.44,0.13}{\textbf{{#1}}}}
\newcommand{\OperatorTok}[1]{\textcolor[rgb]{0.40,0.40,0.40}{{#1}}}
\newcommand{\BuiltInTok}[1]{{#1}}
\newcommand{\ExtensionTok}[1]{{#1}}
\newcommand{\PreprocessorTok}[1]{\textcolor[rgb]{0.74,0.48,0.00}{{#1}}}
\newcommand{\AttributeTok}[1]{\textcolor[rgb]{0.49,0.56,0.16}{{#1}}}
\newcommand{\InformationTok}[1]{\textcolor[rgb]{0.38,0.63,0.69}{\textbf{\textit{{#1}}}}}
\newcommand{\WarningTok}[1]{\textcolor[rgb]{0.38,0.63,0.69}{\textbf{\textit{{#1}}}}}
\makeatletter
\newsavebox\pandoc@box
\newcommand*\pandocbounded[1]{%
	\sbox\pandoc@box{#1}%
	% scaling factors for width and height
	\Gscale@div\@tempa\textheight{\dimexpr\ht\pandoc@box+\dp\pandoc@box\relax}%
	\Gscale@div\@tempb\linewidth{\wd\pandoc@box}%
	% select the smaller of both
	\ifdim\@tempb\p@<\@tempa\p@
	\let\@tempa\@tempb
	\fi
	% scaling accordingly (\@tempa < 1)
	\ifdim\@tempa\p@<\p@
	\scalebox{\@tempa}{\usebox\pandoc@box}%
	% scaling not needed, use as it is
	\else
	\usebox{\pandoc@box}%
	\fi
}
\makeatother

% Define a nice break command that doesn't care if a line doesn't already
% exist.
\def\br{\hspace*{\fill} \\* }
% Math Jax compatibility definitions
\def\gt{>}
\def\lt{<}
\let\Oldtex\TeX
\let\Oldlatex\LaTeX
\renewcommand{\TeX}{\textrm{\Oldtex}}
\renewcommand{\LaTeX}{\textrm{\Oldlatex}}
% Document parameters
% Document title
\title{\textbf{{\huge 机器视觉-实践项目4}}}
\vspace{2em}\author{国嘉通\thanks{Tel.+86-166 8130 6085,\ E-mail:guojiatom2006@outlook.com}} 

% Pygments definitions
\makeatletter
\def\PY@reset{\let\PY@it=\relax \let\PY@bf=\relax%
	\let\PY@ul=\relax \let\PY@tc=\relax%
	\let\PY@bc=\relax \let\PY@ff=\relax}
\def\PY@tok#1{\csname PY@tok@#1\endcsname}
\def\PY@toks#1+{\ifx\relax#1\empty\else%
	\PY@tok{#1}\expandafter\PY@toks\fi}
\def\PY@do#1{\PY@bc{\PY@tc{\PY@ul{%
				\PY@it{\PY@bf{\PY@ff{#1}}}}}}}
\def\PY#1#2{\PY@reset\PY@toks#1+\relax+\PY@do{#2}}

\@namedef{PY@tok@w}{\def\PY@tc##1{\textcolor[rgb]{0.73,0.73,0.73}{##1}}}
\@namedef{PY@tok@c}{\let\PY@it=\textit\def\PY@tc##1{\textcolor[rgb]{0.24,0.48,0.48}{##1}}}
\@namedef{PY@tok@cp}{\def\PY@tc##1{\textcolor[rgb]{0.61,0.40,0.00}{##1}}}
\@namedef{PY@tok@k}{\let\PY@bf=\textbf\def\PY@tc##1{\textcolor[rgb]{0.00,0.50,0.00}{##1}}}
\@namedef{PY@tok@kp}{\def\PY@tc##1{\textcolor[rgb]{0.00,0.50,0.00}{##1}}}
\@namedef{PY@tok@kt}{\def\PY@tc##1{\textcolor[rgb]{0.69,0.00,0.25}{##1}}}
\@namedef{PY@tok@o}{\def\PY@tc##1{\textcolor[rgb]{0.40,0.40,0.40}{##1}}}
\@namedef{PY@tok@ow}{\let\PY@bf=\textbf\def\PY@tc##1{\textcolor[rgb]{0.67,0.13,1.00}{##1}}}
\@namedef{PY@tok@nb}{\def\PY@tc##1{\textcolor[rgb]{0.00,0.50,0.00}{##1}}}
\@namedef{PY@tok@nf}{\def\PY@tc##1{\textcolor[rgb]{0.00,0.00,1.00}{##1}}}
\@namedef{PY@tok@nc}{\let\PY@bf=\textbf\def\PY@tc##1{\textcolor[rgb]{0.00,0.00,1.00}{##1}}}
\@namedef{PY@tok@nn}{\let\PY@bf=\textbf\def\PY@tc##1{\textcolor[rgb]{0.00,0.00,1.00}{##1}}}
\@namedef{PY@tok@ne}{\let\PY@bf=\textbf\def\PY@tc##1{\textcolor[rgb]{0.80,0.25,0.22}{##1}}}
\@namedef{PY@tok@nv}{\def\PY@tc##1{\textcolor[rgb]{0.10,0.09,0.49}{##1}}}
\@namedef{PY@tok@no}{\def\PY@tc##1{\textcolor[rgb]{0.53,0.00,0.00}{##1}}}
\@namedef{PY@tok@nl}{\def\PY@tc##1{\textcolor[rgb]{0.46,0.46,0.00}{##1}}}
\@namedef{PY@tok@ni}{\let\PY@bf=\textbf\def\PY@tc##1{\textcolor[rgb]{0.44,0.44,0.44}{##1}}}
\@namedef{PY@tok@na}{\def\PY@tc##1{\textcolor[rgb]{0.41,0.47,0.13}{##1}}}
\@namedef{PY@tok@nt}{\let\PY@bf=\textbf\def\PY@tc##1{\textcolor[rgb]{0.00,0.50,0.00}{##1}}}
\@namedef{PY@tok@nd}{\def\PY@tc##1{\textcolor[rgb]{0.67,0.13,1.00}{##1}}}
\@namedef{PY@tok@s}{\def\PY@tc##1{\textcolor[rgb]{0.73,0.13,0.13}{##1}}}
\@namedef{PY@tok@sd}{\let\PY@it=\textit\def\PY@tc##1{\textcolor[rgb]{0.73,0.13,0.13}{##1}}}
\@namedef{PY@tok@si}{\let\PY@bf=\textbf\def\PY@tc##1{\textcolor[rgb]{0.64,0.35,0.47}{##1}}}
\@namedef{PY@tok@se}{\let\PY@bf=\textbf\def\PY@tc##1{\textcolor[rgb]{0.67,0.36,0.12}{##1}}}
\@namedef{PY@tok@sr}{\def\PY@tc##1{\textcolor[rgb]{0.64,0.35,0.47}{##1}}}
\@namedef{PY@tok@ss}{\def\PY@tc##1{\textcolor[rgb]{0.10,0.09,0.49}{##1}}}
\@namedef{PY@tok@sx}{\def\PY@tc##1{\textcolor[rgb]{0.00,0.50,0.00}{##1}}}
\@namedef{PY@tok@m}{\def\PY@tc##1{\textcolor[rgb]{0.40,0.40,0.40}{##1}}}
\@namedef{PY@tok@gh}{\let\PY@bf=\textbf\def\PY@tc##1{\textcolor[rgb]{0.00,0.00,0.50}{##1}}}
\@namedef{PY@tok@gu}{\let\PY@bf=\textbf\def\PY@tc##1{\textcolor[rgb]{0.50,0.00,0.50}{##1}}}
\@namedef{PY@tok@gd}{\def\PY@tc##1{\textcolor[rgb]{0.63,0.00,0.00}{##1}}}
\@namedef{PY@tok@gi}{\def\PY@tc##1{\textcolor[rgb]{0.00,0.52,0.00}{##1}}}
\@namedef{PY@tok@gr}{\def\PY@tc##1{\textcolor[rgb]{0.89,0.00,0.00}{##1}}}
\@namedef{PY@tok@ge}{\let\PY@it=\textit}
\@namedef{PY@tok@gs}{\let\PY@bf=\textbf}
\@namedef{PY@tok@ges}{\let\PY@bf=\textbf\let\PY@it=\textit}
\@namedef{PY@tok@gp}{\let\PY@bf=\textbf\def\PY@tc##1{\textcolor[rgb]{0.00,0.00,0.50}{##1}}}
\@namedef{PY@tok@go}{\def\PY@tc##1{\textcolor[rgb]{0.44,0.44,0.44}{##1}}}
\@namedef{PY@tok@gt}{\def\PY@tc##1{\textcolor[rgb]{0.00,0.27,0.87}{##1}}}
\@namedef{PY@tok@err}{\def\PY@bc##1{{\setlength{\fboxsep}{\string -\fboxrule}\fcolorbox[rgb]{1.00,0.00,0.00}{1,1,1}{\strut ##1}}}}
\@namedef{PY@tok@kc}{\let\PY@bf=\textbf\def\PY@tc##1{\textcolor[rgb]{0.00,0.50,0.00}{##1}}}
\@namedef{PY@tok@kd}{\let\PY@bf=\textbf\def\PY@tc##1{\textcolor[rgb]{0.00,0.50,0.00}{##1}}}
\@namedef{PY@tok@kn}{\let\PY@bf=\textbf\def\PY@tc##1{\textcolor[rgb]{0.00,0.50,0.00}{##1}}}
\@namedef{PY@tok@kr}{\let\PY@bf=\textbf\def\PY@tc##1{\textcolor[rgb]{0.00,0.50,0.00}{##1}}}
\@namedef{PY@tok@bp}{\def\PY@tc##1{\textcolor[rgb]{0.00,0.50,0.00}{##1}}}
\@namedef{PY@tok@fm}{\def\PY@tc##1{\textcolor[rgb]{0.00,0.00,1.00}{##1}}}
\@namedef{PY@tok@vc}{\def\PY@tc##1{\textcolor[rgb]{0.10,0.09,0.49}{##1}}}
\@namedef{PY@tok@vg}{\def\PY@tc##1{\textcolor[rgb]{0.10,0.09,0.49}{##1}}}
\@namedef{PY@tok@vi}{\def\PY@tc##1{\textcolor[rgb]{0.10,0.09,0.49}{##1}}}
\@namedef{PY@tok@vm}{\def\PY@tc##1{\textcolor[rgb]{0.10,0.09,0.49}{##1}}}
\@namedef{PY@tok@sa}{\def\PY@tc##1{\textcolor[rgb]{0.73,0.13,0.13}{##1}}}
\@namedef{PY@tok@sb}{\def\PY@tc##1{\textcolor[rgb]{0.73,0.13,0.13}{##1}}}
\@namedef{PY@tok@sc}{\def\PY@tc##1{\textcolor[rgb]{0.73,0.13,0.13}{##1}}}
\@namedef{PY@tok@dl}{\def\PY@tc##1{\textcolor[rgb]{0.73,0.13,0.13}{##1}}}
\@namedef{PY@tok@s2}{\def\PY@tc##1{\textcolor[rgb]{0.73,0.13,0.13}{##1}}}
\@namedef{PY@tok@sh}{\def\PY@tc##1{\textcolor[rgb]{0.73,0.13,0.13}{##1}}}
\@namedef{PY@tok@s1}{\def\PY@tc##1{\textcolor[rgb]{0.73,0.13,0.13}{##1}}}
\@namedef{PY@tok@mb}{\def\PY@tc##1{\textcolor[rgb]{0.40,0.40,0.40}{##1}}}
\@namedef{PY@tok@mf}{\def\PY@tc##1{\textcolor[rgb]{0.40,0.40,0.40}{##1}}}
\@namedef{PY@tok@mh}{\def\PY@tc##1{\textcolor[rgb]{0.40,0.40,0.40}{##1}}}
\@namedef{PY@tok@mi}{\def\PY@tc##1{\textcolor[rgb]{0.40,0.40,0.40}{##1}}}
\@namedef{PY@tok@il}{\def\PY@tc##1{\textcolor[rgb]{0.40,0.40,0.40}{##1}}}
\@namedef{PY@tok@mo}{\def\PY@tc##1{\textcolor[rgb]{0.40,0.40,0.40}{##1}}}
\@namedef{PY@tok@ch}{\let\PY@it=\textit\def\PY@tc##1{\textcolor[rgb]{0.24,0.48,0.48}{##1}}}
\@namedef{PY@tok@cm}{\let\PY@it=\textit\def\PY@tc##1{\textcolor[rgb]{0.24,0.48,0.48}{##1}}}
\@namedef{PY@tok@cpf}{\let\PY@it=\textit\def\PY@tc##1{\textcolor[rgb]{0.24,0.48,0.48}{##1}}}
\@namedef{PY@tok@c1}{\let\PY@it=\textit\def\PY@tc##1{\textcolor[rgb]{0.24,0.48,0.48}{##1}}}
\@namedef{PY@tok@cs}{\let\PY@it=\textit\def\PY@tc##1{\textcolor[rgb]{0.24,0.48,0.48}{##1}}}

\def\PYZbs{\char`\\}
\def\PYZus{\char`\_}
\def\PYZob{\char`\{}
\def\PYZcb{\char`\}}
\def\PYZca{\char`\^}
\def\PYZam{\char`\&}
\def\PYZlt{\char`\<}
\def\PYZgt{\char`\>}
\def\PYZsh{\char`\#}
\def\PYZpc{\char`\%}
\def\PYZdl{\char`\$}
\def\PYZhy{\char`\-}
\def\PYZsq{\char`\'}
\def\PYZdq{\char`\"}
\def\PYZti{\char`\~}
% for compatibility with earlier versions
\def\PYZat{@}
\def\PYZlb{[}
\def\PYZrb{]}
\makeatother


% For linebreaks inside Verbatim environment from package fancyvrb.
\makeatletter
\newbox\Wrappedcontinuationbox
\newbox\Wrappedvisiblespacebox
\newcommand*\Wrappedvisiblespace {\textcolor{red}{\textvisiblespace}}
\newcommand*\Wrappedcontinuationsymbol {\textcolor{red}{\llap{\tiny$\m@th\hookrightarrow$}}}
\newcommand*\Wrappedcontinuationindent {3ex }
\newcommand*\Wrappedafterbreak {\kern\Wrappedcontinuationindent\copy\Wrappedcontinuationbox}
% Take advantage of the already applied Pygments mark-up to insert
% potential linebreaks for TeX processing.
%        {, <, #, %, $, ' and ": go to next line.
	%        _, }, ^, &, >, - and ~: stay at end of broken line.
% Use of \textquotesingle for straight quote.
\newcommand*\Wrappedbreaksatspecials {%
	\def\PYGZus{\discretionary{\char`\_}{\Wrappedafterbreak}{\char`\_}}%
	\def\PYGZob{\discretionary{}{\Wrappedafterbreak\char`\{}{\char`\{}}%
	\def\PYGZcb{\discretionary{\char`\}}{\Wrappedafterbreak}{\char`\}}}%
	\def\PYGZca{\discretionary{\char`\^}{\Wrappedafterbreak}{\char`\^}}%
	\def\PYGZam{\discretionary{\char`\&}{\Wrappedafterbreak}{\char`\&}}%
	\def\PYGZlt{\discretionary{}{\Wrappedafterbreak\char`\<}{\char`\<}}%
	\def\PYGZgt{\discretionary{\char`\>}{\Wrappedafterbreak}{\char`\>}}%
	\def\PYGZsh{\discretionary{}{\Wrappedafterbreak\char`\#}{\char`\#}}%
	\def\PYGZpc{\discretionary{}{\Wrappedafterbreak\char`\%}{\char`\%}}%
	\def\PYGZdl{\discretionary{}{\Wrappedafterbreak\char`\$}{\char`\$}}%
	\def\PYGZhy{\discretionary{\char`\-}{\Wrappedafterbreak}{\char`\-}}%
	\def\PYGZsq{\discretionary{}{\Wrappedafterbreak\textquotesingle}{\textquotesingle}}%
	\def\PYGZdq{\discretionary{}{\Wrappedafterbreak\char`\"}{\char`\"}}%
	\def\PYGZti{\discretionary{\char`\~}{\Wrappedafterbreak}{\char`\~}}%
}
% Some characters . , ; ? ! / are not pygmentized.
% This macro makes them "active" and they will insert potential linebreaks
\newcommand*\Wrappedbreaksatpunct {%
	\lccode`\~`\.\lowercase{\def~}{\discretionary{\hbox{\char`\.}}{\Wrappedafterbreak}{\hbox{\char`\.}}}%
	\lccode`\~`\,\lowercase{\def~}{\discretionary{\hbox{\char`\,}}{\Wrappedafterbreak}{\hbox{\char`\,}}}%
	\lccode`\~`\;\lowercase{\def~}{\discretionary{\hbox{\char`\;}}{\Wrappedafterbreak}{\hbox{\char`\;}}}%
	\lccode`\~`\:\lowercase{\def~}{\discretionary{\hbox{\char`\:}}{\Wrappedafterbreak}{\hbox{\char`\:}}}%
	\lccode`\~`\?\lowercase{\def~}{\discretionary{\hbox{\char`\?}}{\Wrappedafterbreak}{\hbox{\char`\?}}}%
	\lccode`\~`\!\lowercase{\def~}{\discretionary{\hbox{\char`\!}}{\Wrappedafterbreak}{\hbox{\char`\!}}}%
	\lccode`\~`\/\lowercase{\def~}{\discretionary{\hbox{\char`\/}}{\Wrappedafterbreak}{\hbox{\char`\/}}}%
	\catcode`\.\active
	\catcode`\,\active
	\catcode`\;\active
	\catcode`\:\active
	\catcode`\?\active
	\catcode`\!\active
	\catcode`\/\active
	\lccode`\~`\~
}
\makeatother

\let\OriginalVerbatim=\Verbatim
\makeatletter
\renewcommand{\Verbatim}[1][1]{%
	%\parskip\z@skip
	\sbox\Wrappedcontinuationbox {\Wrappedcontinuationsymbol}%
	\sbox\Wrappedvisiblespacebox {\FV@SetupFont\Wrappedvisiblespace}%
	\def\FancyVerbFormatLine ##1{\hsize\linewidth
		\vtop{\raggedright\hyphenpenalty\z@\exhyphenpenalty\z@
			\doublehyphendemerits\z@\finalhyphendemerits\z@
			\strut ##1\strut}%
	}%
	% If the linebreak is at a space, the latter will be displayed as visible
	% space at end of first line, and a continuation symbol starts next line.
	% Stretch/shrink are however usually zero for typewriter font.
	\def\FV@Space {%
		\nobreak\hskip\z@ plus\fontdimen3\font minus\fontdimen4\font
		\discretionary{\copy\Wrappedvisiblespacebox}{\Wrappedafterbreak}
		{\kern\fontdimen2\font}%
	}%
	
	% Allow breaks at special characters using \PYG... macros.
	\Wrappedbreaksatspecials
	% Breaks at punctuation characters . , ; ? ! and / need catcode=\active
	\OriginalVerbatim[#1,codes*=\Wrappedbreaksatpunct]%
}
\makeatother

% Exact colors from NB
\definecolor{incolor}{HTML}{303F9F}
\definecolor{outcolor}{HTML}{D84315}
\definecolor{cellborder}{HTML}{CFCFCF}
\definecolor{cellbackground}{HTML}{F7F7F7}

% prompt
\makeatletter
\newcommand{\boxspacing}{\kern\kvtcb@left@rule\kern\kvtcb@boxsep}
\makeatother
\newcommand{\prompt}[4]{
	{\ttfamily\llap{{\color{#2}[#3]:\hspace{3pt}#4}}\vspace{-\baselineskip}}
}



% Prevent overflowing lines due to hard-to-break entities
\sloppy
% Setup hyperref package
\hypersetup{
	breaklinks=true,  % so long urls are correctly broken across lines
	colorlinks=true,
	urlcolor=urlcolor,
	linkcolor=linkcolor,
	citecolor=citecolor,
}
% Slightly bigger margins than the latex defaults

\geometry{verbose,tmargin=1in,bmargin=1in,lmargin=1in,rmargin=1in}
    
    

\begin{document}
    
    \maketitle
    
      \vspace*{4em}
    \begin{abstract}
    	机器视觉第四次实验课程主要包含两大部分的练习，一是图像形态学处理，二是边缘检测处理。对于图像形态学处理部分，分别从腐蚀、膨胀、开运算、闭运算、骨架提取五个方面进行详细阐释。对于边缘检测处理部分，分别从Sobel 算子与 Canny 算子两个角度进行阐释。另外，还介绍了K-means与Mean Shift两种图像分割方法。希望本篇文章是一篇内容翔实、方便随时翻阅复习学习OpenCV的学习笔记。
    \end{abstract}
    
    \begin{center}
    	\vspace*{5em}{\large \textbf{Now the CODE is available at:}}\\\underline{https://github.com/GplasT0810/Machine-Vision-Experiment.git}\\
    	\vspace*{6em}{\Large \textbf{长沙理工大学}}\\
    	\vspace{0.7em}{\large \textbf{卓越工程师学院\ 绿色智慧交通 2401班}}\\
    	\vspace{0.7em}\textbf{{\large 202404080608}}\\
    	\vspace{0.7em}{\large \textbf{国嘉通}}
    \end{center}
    \newpage
    
    \tableofcontents
    \begin{center}
    	\vspace{11em}\adjustimage{max size={0.9\linewidth}{0.9\paperheight}}{SchPic.png}
    \end{center}
    
    \vspace*{13em}--------------------------------------------\\
    Complated Time: Oct.$1^{st}$.2025\\ \\
    Guo Jiatong , Green \& Intelligent Transportation 2401\\
    Elite Engineering School , Changsha University of Science \& Technology.\\
    Changsha 410114 , Hunan , China
    
    \newpage

    
    \section{形态学处理}
    本部分分别从腐蚀、膨胀、开运算、闭运算、骨架提取五个方面进行详细阐释。
    \subsection{腐蚀运算(Erosion)}
    \begin{itemize}
    	\item 原理\\
    	使用结构元素逐个像素地扫描图像，仅当结构元素完全覆盖前景区域时，图像的中心像素才会被保留。若设图像集合为$A$，结构元素为$B$，则腐蚀运算定义为
    	$$A\ominus B = \{z\mid B_z\subseteq A\}$$
    	即结构元素$B$在图像$A$上滑动，只有当$B$完全包含于$A$时，中心像素才会被保留。
    	\item 作用\\
    	去除小噪点（噪点需要小于结构元素）、分离粘连物体并进行边缘提取前的预处理、缩小前景区域等。
    	\item 代码实现
    	\begin{center}
    			\fbox{
    			\begin{minipage}{14cm}
    				kernel = cv2.getStructuringElement(cv2.MORPH\_RECT, (5, 5))\\
    				eroded = cv2.erode(binary, kernel, iterations=1)
    		\end{minipage}}
    	\end{center}
    \end{itemize}
    \begin{tcolorbox}[breakable, size=fbox, boxrule=1pt, pad at break*=1mm,colback=cellbackground, colframe=cellborder]
%_{}\prompt{In}{incolor}{244}{\boxspacing}
\begin{Verbatim}[commandchars=\\\{\}]
\PY{c+c1}{\PYZsh{}图像腐蚀运算}
\PY{k+kn}{import}\PY{+w}{ }\PY{n+nn}{cv2}
\PY{k+kn}{import}\PY{+w}{ }\PY{n+nn}{numpy}\PY{+w}{ }\PY{k}{as}\PY{+w}{ }\PY{n+nn}{np}
\PY{k+kn}{import}\PY{+w}{ }\PY{n+nn}{matplotlib}\PY{n+nn}{.}\PY{n+nn}{pyplot}\PY{+w}{ }\PY{k}{as}\PY{+w}{ }\PY{n+nn}{plt}

\PY{n}{img\PYZus{}path} \PY{o}{=} \PY{l+s+s1}{\PYZsq{}}\PY{l+s+s1}{/Users/guo2006/myenv/Machine Vision Experiment/Machine Vision Experiment 4/zb.jpg}\PY{l+s+s1}{\PYZsq{}}
\PY{n}{img} \PY{o}{=} \PY{n}{cv2}\PY{o}{.}\PY{n}{imread}\PY{p}{(}\PY{n}{img\PYZus{}path}\PY{p}{)}
\PY{k}{if} \PY{n}{img} \PY{o+ow}{is} \PY{k+kc}{None}\PY{p}{:}
    \PY{k}{raise} \PY{n+ne}{FileNotFoundError}\PY{p}{(}\PY{l+s+s1}{\PYZsq{}}\PY{l+s+s1}{图片路径错误或文件不存在}\PY{l+s+s1}{\PYZsq{}}\PY{p}{)}

\PY{n}{gray} \PY{o}{=} \PY{n}{cv2}\PY{o}{.}\PY{n}{cvtColor}\PY{p}{(}\PY{n}{img}\PY{p}{,} \PY{n}{cv2}\PY{o}{.}\PY{n}{COLOR\PYZus{}BGR2GRAY}\PY{p}{)}
\PY{n}{\PYZus{}}\PY{p}{,} \PY{n}{binary} \PY{o}{=} \PY{n}{cv2}\PY{o}{.}\PY{n}{threshold}\PY{p}{(}\PY{n}{gray}\PY{p}{,} \PY{l+m+mi}{0}\PY{p}{,} \PY{l+m+mi}{255}\PY{p}{,} \PY{n}{cv2}\PY{o}{.}\PY{n}{THRESH\PYZus{}BINARY} \PY{o}{+} \PY{n}{cv2}\PY{o}{.}\PY{n}{THRESH\PYZus{}OTSU}\PY{p}{)} 

\PY{c+c1}{\PYZsh{}定义腐蚀核}
\PY{n}{kernel} \PY{o}{=} \PY{n}{cv2}\PY{o}{.}\PY{n}{getStructuringElement}\PY{p}{(}\PY{n}{cv2}\PY{o}{.}\PY{n}{MORPH\PYZus{}RECT}\PY{p}{,} \PY{p}{(}\PY{l+m+mi}{5}\PY{p}{,} \PY{l+m+mi}{5}\PY{p}{)}\PY{p}{)}
\PY{n}{eroded} \PY{o}{=} \PY{n}{cv2}\PY{o}{.}\PY{n}{erode}\PY{p}{(}\PY{n}{binary}\PY{p}{,} \PY{n}{kernel}\PY{p}{,} \PY{n}{iterations}\PY{o}{=}\PY{l+m+mi}{1}\PY{p}{)}

\PY{n}{vis} \PY{o}{=} \PY{n}{np}\PY{o}{.}\PY{n}{hstack}\PY{p}{(}\PY{p}{[}\PY{n}{gray}\PY{p}{,} \PY{n}{binary}\PY{p}{,} \PY{n}{eroded}\PY{p}{]}\PY{p}{)} 
\PY{n}{vis} \PY{o}{=} \PY{n}{cv2}\PY{o}{.}\PY{n}{cvtColor}\PY{p}{(}\PY{n}{vis}\PY{p}{,} \PY{n}{cv2}\PY{o}{.}\PY{n}{COLOR\PYZus{}BGR2RGB}\PY{p}{)}
\PY{n}{plt}\PY{o}{.}\PY{n}{imshow}\PY{p}{(}\PY{n}{vis}\PY{p}{)}
\PY{n}{plt}\PY{o}{.}\PY{n}{axis}\PY{p}{(}\PY{l+s+s1}{\PYZsq{}}\PY{l+s+s1}{off}\PY{l+s+s1}{\PYZsq{}}\PY{p}{)}
\PY{n}{plt}\PY{o}{.}\PY{n}{title}\PY{p}{(}\PY{l+s+s1}{\PYZsq{}}\PY{l+s+s1}{Gray / Binary / Eroded}\PY{l+s+s1}{\PYZsq{}}\PY{p}{)}
\end{Verbatim}
\end{tcolorbox}

            \begin{tcolorbox}[breakable, size=fbox, boxrule=.5pt, pad at break*=1mm, opacityfill=0]
%\prompt{Out}{outcolor}{244}{\boxspacing}
\begin{Verbatim}[commandchars=\\\{\}]
Text(0.5, 1.0, 'Gray / Binary / Eroded')
\end{Verbatim}
\end{tcolorbox}
        
    \begin{center}
    \adjustimage{max size={0.9\linewidth}{0.9\paperheight}}{output_0_1.png}
    \end{center}
   % { \hspace*{\fill} \\}
    
    
    \subsection{膨胀运算(Dilation)}
    \begin{itemize}
    	\item 原理\\
    	结构元素覆盖任意一个前景像素（结构元素与前景有交集），则中心即设为前景。若设图像集合为$A$，结构元素为$B$，则膨胀运算定义为
    	$$A\oplus B = \{z\mid B_z\cap A\neq\varnothing\}$$
    	\item 作用\\
    	填补小孔洞，连接断裂区域、扩大前景区域等，还可以提取图像边缘，膨胀运算是腐蚀运算的对偶操作。
    	\item 代码实现
    	\begin{center}
    		\fbox{
    			\begin{minipage}{14cm}
    				kernel = cv2.getStructuringElement(cv2.MORPH\_RECT, (5, 5))\\
    				dilated = cv2.dilate(binary, kernel, iterations=1)
    		\end{minipage}}
    	\end{center}
    \end{itemize}
    
    \begin{tcolorbox}[breakable, size=fbox, boxrule=1pt, pad at break*=1mm,colback=cellbackground, colframe=cellborder]
%\prompt{In}{incolor}{245}{\boxspacing}
\begin{Verbatim}[commandchars=\\\{\}]
\PY{c+c1}{\PYZsh{}图像膨胀运算}
\PY{k+kn}{import}\PY{+w}{ }\PY{n+nn}{cv2}
\PY{k+kn}{import}\PY{+w}{ }\PY{n+nn}{numpy}\PY{+w}{ }\PY{k}{as}\PY{+w}{ }\PY{n+nn}{np}
\PY{k+kn}{import}\PY{+w}{ }\PY{n+nn}{matplotlib}\PY{n+nn}{.}\PY{n+nn}{pyplot}\PY{+w}{ }\PY{k}{as}\PY{+w}{ }\PY{n+nn}{plt}

\PY{n}{img\PYZus{}path} \PY{o}{=} \PY{l+s+s1}{\PYZsq{}}\PY{l+s+s1}{/Users/guo2006/myenv/Machine Vision Experiment/Machine Vision Experiment 4/zb.jpg}\PY{l+s+s1}{\PYZsq{}}
\PY{n}{img} \PY{o}{=} \PY{n}{cv2}\PY{o}{.}\PY{n}{imread}\PY{p}{(}\PY{n}{img\PYZus{}path}\PY{p}{)}
\PY{k}{if} \PY{n}{img} \PY{o+ow}{is} \PY{k+kc}{None}\PY{p}{:}
    \PY{k}{raise} \PY{n+ne}{FileNotFoundError}\PY{p}{(}\PY{l+s+s1}{\PYZsq{}}\PY{l+s+s1}{图片路径错误或文件不存在}\PY{l+s+s1}{\PYZsq{}}\PY{p}{)}

\PY{c+c1}{\PYZsh{}图像灰度处理与二值化}
\PY{n}{gray} \PY{o}{=} \PY{n}{cv2}\PY{o}{.}\PY{n}{cvtColor}\PY{p}{(}\PY{n}{img}\PY{p}{,} \PY{n}{cv2}\PY{o}{.}\PY{n}{COLOR\PYZus{}BGR2GRAY}\PY{p}{)}
\PY{n}{\PYZus{}}\PY{p}{,} \PY{n}{binary} \PY{o}{=} \PY{n}{cv2}\PY{o}{.}\PY{n}{threshold}\PY{p}{(}\PY{n}{gray}\PY{p}{,} \PY{l+m+mi}{0}\PY{p}{,} \PY{l+m+mi}{255}\PY{p}{,}\PY{n}{cv2}\PY{o}{.}\PY{n}{THRESH\PYZus{}BINARY} \PY{o}{+} \PY{n}{cv2}\PY{o}{.}\PY{n}{THRESH\PYZus{}OTSU}\PY{p}{)}
\PY{c+c1}{\PYZsh{}前景为黑色\PYZhy{}\PYZgt{}THRESH\PYZus{}BINARY}
\PY{c+c1}{\PYZsh{}前景为白色\PYZhy{}\PYZgt{}THRESH\PYZus{}BINARY\PYZus{}INV}

\PY{c+c1}{\PYZsh{}定义膨胀核}
\PY{n}{kernel} \PY{o}{=} \PY{n}{cv2}\PY{o}{.}\PY{n}{getStructuringElement}\PY{p}{(}\PY{n}{cv2}\PY{o}{.}\PY{n}{MORPH\PYZus{}RECT}\PY{p}{,} \PY{p}{(}\PY{l+m+mi}{5}\PY{p}{,} \PY{l+m+mi}{5}\PY{p}{)}\PY{p}{)}
\PY{n}{dilated} \PY{o}{=} \PY{n}{cv2}\PY{o}{.}\PY{n}{dilate}\PY{p}{(}\PY{n}{binary}\PY{p}{,} \PY{n}{kernel}\PY{p}{,} \PY{n}{iterations}\PY{o}{=}\PY{l+m+mi}{1}\PY{p}{)}

\PY{n}{vis} \PY{o}{=} \PY{n}{np}\PY{o}{.}\PY{n}{hstack}\PY{p}{(}\PY{p}{[}\PY{n}{gray}\PY{p}{,} \PY{n}{binary}\PY{p}{,} \PY{n}{dilated}\PY{p}{]}\PY{p}{)}
\PY{n}{vis} \PY{o}{=} \PY{n}{cv2}\PY{o}{.}\PY{n}{cvtColor}\PY{p}{(}\PY{n}{vis}\PY{p}{,} \PY{n}{cv2}\PY{o}{.}\PY{n}{COLOR\PYZus{}BGR2RGB}\PY{p}{)}
\PY{n}{plt}\PY{o}{.}\PY{n}{imshow}\PY{p}{(}\PY{n}{vis}\PY{p}{)}
\PY{n}{plt}\PY{o}{.}\PY{n}{axis}\PY{p}{(}\PY{l+s+s1}{\PYZsq{}}\PY{l+s+s1}{off}\PY{l+s+s1}{\PYZsq{}}\PY{p}{)}
\PY{n}{plt}\PY{o}{.}\PY{n}{title}\PY{p}{(}\PY{l+s+s1}{\PYZsq{}}\PY{l+s+s1}{Gray / Binary / Dilated}\PY{l+s+s1}{\PYZsq{}}\PY{p}{)}
\end{Verbatim}
\end{tcolorbox}



            \begin{tcolorbox}[breakable, size=fbox, boxrule=.5pt, pad at break*=1mm, opacityfill=0]
%\prompt{Out}{outcolor}{245}{\boxspacing}
\begin{Verbatim}[commandchars=\\\{\}]
Text(0.5, 1.0, 'Gray / Binary / Dilated')
\end{Verbatim}
\end{tcolorbox}
        
        
    \begin{center}
    \adjustimage{max size={0.9\linewidth}{0.9\paperheight}}{output_1_1.png}
    \end{center}
  %  { \hspace*{\fill} \\}
    
   
   
         \subsection{开运算(Opening)}
   \begin{itemize}
   	\item 定义\\
   	开运算本质上是腐蚀运算与膨胀运算的叠加，即先进行腐蚀运算，再进行膨胀运算。若设图像集合为$A$，结构元素为$B$，则膨胀运算定义为
   	$$\text{Opening}(A)=(A\ominus B)\oplus B$$
   	\item 作用\\
   	其可以去除小物体，同时保留大物体形状不变，还可以用于平滑轮廓，断开狭窄连接，此外，它还能够抑制亮噪声。在实际应用中，适用于文档图像去噪、图像分割的预处理（清理噪点）、提取大结构忽略局部细节等等场景。
   	\item 代码实现
   	\begin{center}
   		\fbox{
   			\begin{minipage}{14cm}
   				kernel = cv2.getStructuringElement(cv2.MORPH\_RECT, (5, 5))\\
   				opening = cv2.morphologyEx(binary, cv2.MORPH\_OPEN, kernel)
   		\end{minipage}}
   	\end{center}
   \end{itemize}
   
    \begin{tcolorbox}[breakable, size=fbox, boxrule=1pt, pad at break*=1mm,colback=cellbackground, colframe=cellborder]
%\prompt{In}{incolor}{246}{\boxspacing}
\begin{Verbatim}[commandchars=\\\{\}]
\PY{c+c1}{\PYZsh{}图像开运算}
\PY{k+kn}{import}\PY{+w}{ }\PY{n+nn}{cv2}
\PY{k+kn}{import}\PY{+w}{ }\PY{n+nn}{numpy}\PY{+w}{ }\PY{k}{as}\PY{+w}{ }\PY{n+nn}{np}
\PY{k+kn}{import}\PY{+w}{ }\PY{n+nn}{matplotlib}\PY{n+nn}{.}\PY{n+nn}{pyplot}\PY{+w}{ }\PY{k}{as}\PY{+w}{ }\PY{n+nn}{plt}

\PY{n}{img\PYZus{}path} \PY{o}{=} \PY{l+s+s1}{\PYZsq{}}\PY{l+s+s1}{/Users/guo2006/myenv/Machine Vision Experiment/Machine Vision Experiment 4/zb.jpg}\PY{l+s+s1}{\PYZsq{}}
\PY{n}{img} \PY{o}{=} \PY{n}{cv2}\PY{o}{.}\PY{n}{imread}\PY{p}{(}\PY{n}{img\PYZus{}path}\PY{p}{)}
\PY{k}{if} \PY{n}{img} \PY{o+ow}{is} \PY{k+kc}{None}\PY{p}{:}
    \PY{k}{raise} \PY{n+ne}{FileNotFoundError}\PY{p}{(}\PY{l+s+s1}{\PYZsq{}}\PY{l+s+s1}{图片路径错误或文件不存在}\PY{l+s+s1}{\PYZsq{}}\PY{p}{)}

\PY{c+c1}{\PYZsh{}灰度变换与二值化}
\PY{n}{gray} \PY{o}{=} \PY{n}{cv2}\PY{o}{.}\PY{n}{cvtColor}\PY{p}{(}\PY{n}{img}\PY{p}{,} \PY{n}{cv2}\PY{o}{.}\PY{n}{COLOR\PYZus{}BGR2GRAY}\PY{p}{)}
\PY{n}{\PYZus{}}\PY{p}{,} \PY{n}{binary} \PY{o}{=} \PY{n}{cv2}\PY{o}{.}\PY{n}{threshold}\PY{p}{(}\PY{n}{gray}\PY{p}{,} \PY{l+m+mi}{0}\PY{p}{,} \PY{l+m+mi}{255}\PY{p}{,} \PY{n}{cv2}\PY{o}{.}\PY{n}{THRESH\PYZus{}BINARY} \PY{o}{+} \PY{n}{cv2}\PY{o}{.}\PY{n}{THRESH\PYZus{}OTSU}\PY{p}{)}

\PY{c+c1}{\PYZsh{}定义结构元}
\PY{n}{kernel} \PY{o}{=} \PY{n}{cv2}\PY{o}{.}\PY{n}{getStructuringElement}\PY{p}{(}\PY{n}{cv2}\PY{o}{.}\PY{n}{MORPH\PYZus{}RECT}\PY{p}{,} \PY{p}{(}\PY{l+m+mi}{5}\PY{p}{,} \PY{l+m+mi}{5}\PY{p}{)}\PY{p}{)}
\PY{n}{opening} \PY{o}{=} \PY{n}{cv2}\PY{o}{.}\PY{n}{morphologyEx}\PY{p}{(}\PY{n}{binary}\PY{p}{,} \PY{n}{cv2}\PY{o}{.}\PY{n}{MORPH\PYZus{}OPEN}\PY{p}{,} \PY{n}{kernel}\PY{p}{,} \PY{n}{iterations}\PY{o}{=}\PY{l+m+mi}{1}\PY{p}{)}

\PY{n}{vis} \PY{o}{=} \PY{n}{np}\PY{o}{.}\PY{n}{hstack}\PY{p}{(}\PY{p}{[}\PY{n}{gray}\PY{p}{,} \PY{n}{binary}\PY{p}{,} \PY{n}{opening}\PY{p}{]}\PY{p}{)}
\PY{n}{vis} \PY{o}{=} \PY{n}{cv2}\PY{o}{.}\PY{n}{cvtColor}\PY{p}{(}\PY{n}{vis}\PY{p}{,} \PY{n}{cv2}\PY{o}{.}\PY{n}{COLOR\PYZus{}BGR2RGB}\PY{p}{)}
\PY{n}{plt}\PY{o}{.}\PY{n}{imshow}\PY{p}{(}\PY{n}{vis}\PY{p}{)}
\PY{n}{plt}\PY{o}{.}\PY{n}{title}\PY{p}{(}\PY{l+s+s1}{\PYZsq{}}\PY{l+s+s1}{Gray / Binary / Opening}\PY{l+s+s1}{\PYZsq{}}\PY{p}{)}
\PY{n}{plt}\PY{o}{.}\PY{n}{axis}\PY{p}{(}\PY{l+s+s1}{\PYZsq{}}\PY{l+s+s1}{off}\PY{l+s+s1}{\PYZsq{}}\PY{p}{)}
\end{Verbatim}
\end{tcolorbox}



            \begin{tcolorbox}[breakable, size=fbox, boxrule=.5pt, pad at break*=1mm, opacityfill=0]
%\prompt{Out}{outcolor}{246}{\boxspacing}
\begin{Verbatim}[commandchars=\\\{\}]
(np.float64(-0.5), np.float64(3071.5), np.float64(1407.5), np.float64(-0.5))
\end{Verbatim}
\end{tcolorbox}
        
        
    \begin{center}
    \adjustimage{max size={0.9\linewidth}{0.9\paperheight}}{output_2_1.png}
    \end{center}
    %{ \hspace*{\fill} \\}
    
    
 \subsection{闭运算(Closing)}
\begin{itemize}
	\item 定义\\
	闭运算本质上也是腐蚀运算与膨胀运算的叠加，只不过其顺序恰好与开运算相反。其是先进行膨胀运算，再进行腐蚀运算。若设图像集合为$A$，结构元素为$B$，则膨胀运算定义为
	$$\text{Closing}(A)=(A\oplus B)\ominus B$$
	\item 作用\\
	填补前景内部的小空洞，连接断裂部分，平滑边界的同时能消除暗噪声，其有保持物体面积基本不变的优良特性。因此其适合用于连接断裂的字符或线条以及在医学图像处理领域中的组织区域修复等。
	\item 代码实现
	\begin{center}
		\fbox{
			\begin{minipage}{14cm}
				kernel = cv2.getStructuringElement(cv2.MORPH\_RECT, (5, 5)\\
				closing = cv2.morphologyEx(binary, cv2.MORPH\_CLOSE, kernel)
		\end{minipage}}
	\end{center}
\end{itemize}

    \begin{tcolorbox}[breakable, size=fbox, boxrule=1pt, pad at break*=1mm,colback=cellbackground, colframe=cellborder]
%\prompt{In}{incolor}{247}{\boxspacing}
\begin{Verbatim}[commandchars=\\\{\}]
\PY{c+c1}{\PYZsh{}图像闭运算}
\PY{k+kn}{import}\PY{+w}{ }\PY{n+nn}{cv2}
\PY{k+kn}{import}\PY{+w}{ }\PY{n+nn}{numpy}\PY{+w}{ }\PY{k}{as}\PY{+w}{ }\PY{n+nn}{np}
\PY{k+kn}{import}\PY{+w}{ }\PY{n+nn}{matplotlib}\PY{n+nn}{.}\PY{n+nn}{pyplot}\PY{+w}{ }\PY{k}{as}\PY{+w}{ }\PY{n+nn}{plt}

\PY{n}{img\PYZus{}path} \PY{o}{=} \PY{l+s+s1}{\PYZsq{}}\PY{l+s+s1}{/Users/guo2006/myenv/Machine Vision Experiment/Machine Vision Experiment 4/zb.jpg}\PY{l+s+s1}{\PYZsq{}}
\PY{n}{img} \PY{o}{=} \PY{n}{cv2}\PY{o}{.}\PY{n}{imread}\PY{p}{(}\PY{n}{img\PYZus{}path}\PY{p}{)}
\PY{k}{if} \PY{n}{img} \PY{o+ow}{is} \PY{k+kc}{None}\PY{p}{:}
    \PY{k}{raise} \PY{n+ne}{FileNotFoundError}\PY{p}{(}\PY{l+s+s1}{\PYZsq{}}\PY{l+s+s1}{图片路径错误或文件不存在}\PY{l+s+s1}{\PYZsq{}}\PY{p}{)}

\PY{c+c1}{\PYZsh{}灰度变换与二值化}
\PY{n}{gray} \PY{o}{=} \PY{n}{cv2}\PY{o}{.}\PY{n}{cvtColor}\PY{p}{(}\PY{n}{img}\PY{p}{,} \PY{n}{cv2}\PY{o}{.}\PY{n}{COLOR\PYZus{}BGR2GRAY}\PY{p}{)}
\PY{n}{\PYZus{}}\PY{p}{,} \PY{n}{binary} \PY{o}{=} \PY{n}{cv2}\PY{o}{.}\PY{n}{threshold}\PY{p}{(}\PY{n}{gray}\PY{p}{,} \PY{l+m+mi}{0}\PY{p}{,} \PY{l+m+mi}{255}\PY{p}{,} \PY{n}{cv2}\PY{o}{.}\PY{n}{THRESH\PYZus{}BINARY} \PY{o}{+} \PY{n}{cv2}\PY{o}{.}\PY{n}{THRESH\PYZus{}OTSU}\PY{p}{)}

\PY{c+c1}{\PYZsh{}定义结构元}
\PY{n}{kernel} \PY{o}{=} \PY{n}{cv2}\PY{o}{.}\PY{n}{getStructuringElement}\PY{p}{(}\PY{n}{cv2}\PY{o}{.}\PY{n}{MORPH\PYZus{}RECT}\PY{p}{,} \PY{p}{(}\PY{l+m+mi}{5}\PY{p}{,} \PY{l+m+mi}{5}\PY{p}{)}\PY{p}{)}
\PY{n}{closing} \PY{o}{=} \PY{n}{cv2}\PY{o}{.}\PY{n}{morphologyEx}\PY{p}{(}\PY{n}{binary}\PY{p}{,} \PY{n}{cv2}\PY{o}{.}\PY{n}{MORPH\PYZus{}CLOSE}\PY{p}{,} \PY{n}{kernel}\PY{p}{,} \PY{n}{iterations}\PY{o}{=}\PY{l+m+mi}{1}\PY{p}{)}

\PY{n}{vis} \PY{o}{=} \PY{n}{np}\PY{o}{.}\PY{n}{hstack}\PY{p}{(}\PY{p}{[}\PY{n}{gray}\PY{p}{,} \PY{n}{binary}\PY{p}{,} \PY{n}{closing}\PY{p}{]}\PY{p}{)}
\PY{n}{vis} \PY{o}{=} \PY{n}{cv2}\PY{o}{.}\PY{n}{cvtColor}\PY{p}{(}\PY{n}{vis}\PY{p}{,} \PY{n}{cv2}\PY{o}{.}\PY{n}{COLOR\PYZus{}BGR2RGB}\PY{p}{)}
\PY{n}{plt}\PY{o}{.}\PY{n}{imshow}\PY{p}{(}\PY{n}{vis}\PY{p}{)}
\PY{n}{plt}\PY{o}{.}\PY{n}{title}\PY{p}{(}\PY{l+s+s1}{\PYZsq{}}\PY{l+s+s1}{Gray / Binary / Closing}\PY{l+s+s1}{\PYZsq{}}\PY{p}{)}
\PY{n}{plt}\PY{o}{.}\PY{n}{axis}\PY{p}{(}\PY{l+s+s1}{\PYZsq{}}\PY{l+s+s1}{off}\PY{l+s+s1}{\PYZsq{}}\PY{p}{)}
\end{Verbatim}
\end{tcolorbox}

            \begin{tcolorbox}[breakable, size=fbox, boxrule=.5pt, pad at break*=1mm, opacityfill=0]
%\prompt{Out}{outcolor}{247}{\boxspacing}
\begin{Verbatim}[commandchars=\\\{\}]
(np.float64(-0.5), np.float64(3071.5), np.float64(1407.5), np.float64(-0.5))
\end{Verbatim}
\end{tcolorbox}
        
    \begin{center}
    \adjustimage{max size={0.9\linewidth}{0.9\paperheight}}{output_3_1.png}
    \end{center}
  %  { \hspace*{\fill} \\}
    
    
      \subsection{骨架提取(Skeletonization)}
    \begin{itemize}
    	\item 原理\\
    	骨架是前景区域的“中轴线”，骨架提取操作能将前景区域细化为单像素宽的“骨架”，同时保持其拓扑结构不变，thinning至单像素宽。具体算法的实现基于Zhang-Suen 快速并行细化算法、Guo-Hall 算法以及基于距离变换的方法等。
    	\item 应用\\
    	由于其具有骨架单像素宽、能够很好地保持图像的连通性等优良特性，其经常会被用于字符识别、形状分析、路径提取与规划等领域。具体来说，在手写数字或字母的识别，血管、道路裂纹的提取以及图像压缩与结构分析等领域均有应用。
    	\item 代码实现
    	\begin{center}
    		\fbox{
    			\begin{minipage}{14cm}
    				from skimage.morphology import skeletonize\\
    				skeleton = skeletonize(binary\_bool)
    		\end{minipage}}
    	\end{center}
    \end{itemize}
    
    \begin{tcolorbox}[breakable, size=fbox, boxrule=1pt, pad at break*=1mm,colback=cellbackground, colframe=cellborder]
%\prompt{In}{incolor}{248}{\boxspacing}
\begin{Verbatim}[commandchars=\\\{\}]
\PY{c+c1}{\PYZsh{} 骨架操作}
\PY{k+kn}{import}\PY{+w}{ }\PY{n+nn}{cv2}
\PY{k+kn}{import}\PY{+w}{ }\PY{n+nn}{numpy}\PY{+w}{ }\PY{k}{as}\PY{+w}{ }\PY{n+nn}{np}
\PY{k+kn}{import}\PY{+w}{ }\PY{n+nn}{matplotlib}\PY{n+nn}{.}\PY{n+nn}{pyplot}\PY{+w}{ }\PY{k}{as}\PY{+w}{ }\PY{n+nn}{plt}
\PY{k+kn}{from}\PY{+w}{ }\PY{n+nn}{skimage}\PY{n+nn}{.}\PY{n+nn}{morphology}\PY{+w}{ }\PY{k+kn}{import} \PY{n}{skeletonize}
\PY{k+kn}{from}\PY{+w}{ }\PY{n+nn}{skimage}\PY{n+nn}{.}\PY{n+nn}{util}\PY{+w}{ }\PY{k+kn}{import} \PY{n}{img\PYZus{}as\PYZus{}bool}
\PY{c+c1}{\PYZsh{}OpenCV没有单步骨架API，采用Zhang\PYZhy{}Suen快速并行细化算法, 安装scikit\PYZhy{}image}

\PY{n}{img\PYZus{}path} \PY{o}{=} \PY{l+s+s1}{\PYZsq{}}\PY{l+s+s1}{/Users/guo2006/myenv/Machine Vision Experiment/Machine Vision Experiment 4/zb.jpg}\PY{l+s+s1}{\PYZsq{}}
\PY{n}{img} \PY{o}{=} \PY{n}{cv2}\PY{o}{.}\PY{n}{imread}\PY{p}{(}\PY{n}{img\PYZus{}path}\PY{p}{)}
\PY{k}{if} \PY{n}{img} \PY{o+ow}{is} \PY{k+kc}{None}\PY{p}{:}
    \PY{k}{raise} \PY{n+ne}{FileNotFoundError}\PY{p}{(}\PY{l+s+s1}{\PYZsq{}}\PY{l+s+s1}{图片路径错误或文件不存在}\PY{l+s+s1}{\PYZsq{}}\PY{p}{)}

\PY{c+c1}{\PYZsh{}灰度转换与二值化}
\PY{n}{gray} \PY{o}{=} \PY{n}{cv2}\PY{o}{.}\PY{n}{cvtColor}\PY{p}{(}\PY{n}{img}\PY{p}{,} \PY{n}{cv2}\PY{o}{.}\PY{n}{COLOR\PYZus{}BGR2GRAY}\PY{p}{)}
\PY{n}{\PYZus{}}\PY{p}{,} \PY{n}{binary} \PY{o}{=} \PY{n}{cv2}\PY{o}{.}\PY{n}{threshold}\PY{p}{(}\PY{n}{gray}\PY{p}{,} \PY{l+m+mi}{0}\PY{p}{,} \PY{l+m+mi}{255}\PY{p}{,} \PY{n}{cv2}\PY{o}{.}\PY{n}{THRESH\PYZus{}BINARY\PYZus{}INV} \PY{o}{+} \PY{n}{cv2}\PY{o}{.}\PY{n}{THRESH\PYZus{}OTSU}\PY{p}{)}

\PY{n}{binary\PYZus{}bool} \PY{o}{=} \PY{n}{img\PYZus{}as\PYZus{}bool}\PY{p}{(}\PY{n}{binary} \PY{o}{/}\PY{o}{/} \PY{l+m+mi}{255}\PY{p}{)}
\PY{n}{skeleton} \PY{o}{=} \PY{n}{skeletonize}\PY{p}{(}\PY{n}{binary\PYZus{}bool}\PY{p}{)}

\PY{n}{skeleton\PYZus{}uint8} \PY{o}{=} \PY{n}{skeleton}\PY{o}{.}\PY{n}{astype}\PY{p}{(}\PY{n}{np}\PY{o}{.}\PY{n}{uint8}\PY{p}{)} \PY{o}{*} \PY{l+m+mi}{255}
\PY{n}{kernel} \PY{o}{=} \PY{n}{cv2}\PY{o}{.}\PY{n}{getStructuringElement}\PY{p}{(}\PY{n}{cv2}\PY{o}{.}\PY{n}{MORPH\PYZus{}CROSS}\PY{p}{,} \PY{p}{(}\PY{l+m+mi}{3}\PY{p}{,}\PY{l+m+mi}{3}\PY{p}{)}\PY{p}{)}
\PY{n}{skeleton\PYZus{}thick} \PY{o}{=} \PY{n}{cv2}\PY{o}{.}\PY{n}{dilate}\PY{p}{(}\PY{n}{skeleton\PYZus{}uint8}\PY{p}{,} \PY{n}{kernel}\PY{p}{,} \PY{n}{iterations}\PY{o}{=}\PY{l+m+mi}{1}\PY{p}{)}

\PY{n}{vis} \PY{o}{=} \PY{n}{np}\PY{o}{.}\PY{n}{hstack}\PY{p}{(}\PY{p}{[}\PY{n}{gray}\PY{p}{,} \PY{n}{binary}\PY{p}{,} \PY{n}{skeleton\PYZus{}thick}\PY{p}{]}\PY{p}{)}
\PY{n}{plt}\PY{o}{.}\PY{n}{figure}\PY{p}{(}\PY{n}{figsize}\PY{o}{=}\PY{p}{(}\PY{l+m+mi}{12}\PY{p}{,}\PY{l+m+mi}{4}\PY{p}{)}\PY{p}{)}
\PY{n}{plt}\PY{o}{.}\PY{n}{imshow}\PY{p}{(}\PY{n}{vis}\PY{p}{,} \PY{n}{cmap}\PY{o}{=}\PY{l+s+s1}{\PYZsq{}}\PY{l+s+s1}{gray}\PY{l+s+s1}{\PYZsq{}}\PY{p}{)}
\PY{n}{plt}\PY{o}{.}\PY{n}{title}\PY{p}{(}\PY{l+s+s1}{\PYZsq{}}\PY{l+s+s1}{Gray / Binary / Skeleton}\PY{l+s+s1}{\PYZsq{}}\PY{p}{)}
\PY{n}{plt}\PY{o}{.}\PY{n}{axis}\PY{p}{(}\PY{l+s+s1}{\PYZsq{}}\PY{l+s+s1}{off}\PY{l+s+s1}{\PYZsq{}}\PY{p}{)}
\PY{n}{plt}\PY{o}{.}\PY{n}{tight\PYZus{}layout}\PY{p}{(}\PY{p}{)}
\PY{n}{plt}\PY{o}{.}\PY{n}{show}\PY{p}{(}\PY{p}{)}
\end{Verbatim}
\end{tcolorbox}


    \begin{center}
    \adjustimage{max size={0.9\linewidth}{0.9\paperheight}}{output_4_0.png}
    \end{center}
  %  { \hspace*{\fill} \\}
    
    
  \section{边缘检测操作}
  本章节分为基于Sobel算子、Canny算子的边缘检测操作。
  \subsection*{Sobel算子}
  \begin{itemize}
  	\item 原理\\
  	Sobel算子基于一阶导数的近似，分别计算水平（$x$方向）与垂直（$y$方向）梯度，并合成梯度幅值：
  	$$G_x=\begin{bmatrix}-1&0&1\\-2&0&2\\-1&0&1\end{bmatrix},\quad
  	G_y=\begin{bmatrix}-1&-2&-1\\0&0&0\\1&2&1\end{bmatrix},\quad
  	|G|=\sqrt{G_x^2+G_y^2}$$
  	\item 特点\\
  	其具有简单、快速高效的优点，不过其对噪声敏感，且边缘较粗，定位精准性一般，可用于初步的边缘提取或搭配Canny算法共同使用（作为Canny算法的子步骤）。
  	\item 代码实现
  	\begin{center}
  		\fbox{
  		\begin{minipage}{14cm}
  			sobelx = cv2.Sobel(gray, cv2.CV\_16S, 1, 0, ksize=3)\\
  			sobely = cv2.Sobel(gray, cv2.CV\_16S, 0, 1, ksize=3)\\
  			sobel = cv2.magnitude(sobelx.astype(np.float32), sobely.astype(np.float32))\\
  			sobel = cv2.convertScaleAbs(sobel)
  		\end{minipage}}
  	\end{center}
  \end{itemize}
  
      \subsection*{Canny算子}
  \begin{itemize}
  	\item 步骤\\
  	高斯滤波去噪 → 计算梯度幅值与方向 → 非极大值抑制（细化边缘） → 双阈值检测 → 边缘连接（滞后阈值）
  	\item 特点\\
  	边缘细、连续、定位准确且抗噪能力强，但是对参数敏感。建议低阈值为高阈值的$\frac{1}{2}-\frac{1}{3}$。基于以上特点，易于得知Canny算法比Sobel算法更具优势（或者说更具应用前景），其可以被用于目标检测、轮廓提取等，还可以用于图像分割的预处理以及工业缺陷的检测等多个领域。
  	\item 代码实现
  	\begin{center}
  		\fbox{
  			\begin{minipage}{14cm}
  				canny = cv2.Canny(gray, 50, 150)
  		\end{minipage}}
  	\end{center}
  \end{itemize}
    \begin{tcolorbox}[breakable, size=fbox, boxrule=1pt, pad at break*=1mm,colback=cellbackground, colframe=cellborder]
%\prompt{In}{incolor}{249}{\boxspacing}
\begin{Verbatim}[commandchars=\\\{\}]
\PY{c+c1}{\PYZsh{}边缘检测}
\PY{k+kn}{import}\PY{+w}{ }\PY{n+nn}{cv2}
\PY{k+kn}{import}\PY{+w}{ }\PY{n+nn}{numpy}\PY{+w}{ }\PY{k}{as}\PY{+w}{ }\PY{n+nn}{np}
\PY{k+kn}{import}\PY{+w}{ }\PY{n+nn}{matplotlib}\PY{n+nn}{.}\PY{n+nn}{pyplot}\PY{+w}{ }\PY{k}{as}\PY{+w}{ }\PY{n+nn}{plt}

\PY{n}{img\PYZus{}path} \PY{o}{=} \PY{l+s+s1}{\PYZsq{}}\PY{l+s+s1}{/Users/guo2006/myenv/Machine Vision Experiment/Machine Vision Experiment 4/zb.jpg}\PY{l+s+s1}{\PYZsq{}}
\PY{n}{img} \PY{o}{=} \PY{n}{cv2}\PY{o}{.}\PY{n}{imread}\PY{p}{(}\PY{n}{img\PYZus{}path}\PY{p}{)}
\PY{k}{if} \PY{n}{img} \PY{o+ow}{is} \PY{k+kc}{None}\PY{p}{:}
    \PY{k}{raise} \PY{n+ne}{FileNotFoundError}\PY{p}{(}\PY{l+s+s1}{\PYZsq{}}\PY{l+s+s1}{图片路径错误}\PY{l+s+s1}{\PYZsq{}}\PY{p}{)}
\PY{n}{gray} \PY{o}{=} \PY{n}{cv2}\PY{o}{.}\PY{n}{cvtColor}\PY{p}{(}\PY{n}{img}\PY{p}{,} \PY{n}{cv2}\PY{o}{.}\PY{n}{COLOR\PYZus{}BGR2GRAY}\PY{p}{)}

\PY{c+c1}{\PYZsh{}Sobel算子}
\PY{n}{sobelx} \PY{o}{=} \PY{n}{cv2}\PY{o}{.}\PY{n}{Sobel}\PY{p}{(}\PY{n}{gray}\PY{p}{,} \PY{n}{cv2}\PY{o}{.}\PY{n}{CV\PYZus{}16S}\PY{p}{,} \PY{l+m+mi}{1}\PY{p}{,} \PY{l+m+mi}{0}\PY{p}{,} \PY{n}{ksize}\PY{o}{=}\PY{l+m+mi}{3}\PY{p}{)} \PY{c+c1}{\PYZsh{}水平梯度}
\PY{n}{sobely} \PY{o}{=} \PY{n}{cv2}\PY{o}{.}\PY{n}{Sobel}\PY{p}{(}\PY{n}{gray}\PY{p}{,} \PY{n}{cv2}\PY{o}{.}\PY{n}{CV\PYZus{}16S}\PY{p}{,} \PY{l+m+mi}{0}\PY{p}{,} \PY{l+m+mi}{1}\PY{p}{,} \PY{n}{ksize}\PY{o}{=}\PY{l+m+mi}{3}\PY{p}{)} \PY{c+c1}{\PYZsh{}垂直梯度}
\PY{n}{sobel\PYZus{}f} \PY{o}{=} \PY{n}{cv2}\PY{o}{.}\PY{n}{magnitude}\PY{p}{(}\PY{n}{sobelx}\PY{o}{.}\PY{n}{astype}\PY{p}{(}\PY{n}{np}\PY{o}{.}\PY{n}{float32}\PY{p}{)}\PY{p}{,} \PY{n}{sobely}\PY{o}{.}\PY{n}{astype}\PY{p}{(}\PY{n}{np}\PY{o}{.}\PY{n}{float32}\PY{p}{)}\PY{p}{)} \PY{c+c1}{\PYZsh{}合成梯度幅值}
\PY{n}{sobel} \PY{o}{=} \PY{n}{cv2}\PY{o}{.}\PY{n}{convertScaleAbs}\PY{p}{(}\PY{n}{sobel\PYZus{}f}\PY{p}{)} \PY{c+c1}{\PYZsh{}转8位}

\PY{c+c1}{\PYZsh{}Canny算子}
\PY{n}{canny} \PY{o}{=} \PY{n}{cv2}\PY{o}{.}\PY{n}{Canny}\PY{p}{(}\PY{n}{gray}\PY{p}{,} \PY{l+m+mi}{10}\PY{p}{,} \PY{l+m+mi}{100}\PY{p}{)}

\PY{n}{vis} \PY{o}{=} \PY{n}{np}\PY{o}{.}\PY{n}{hstack}\PY{p}{(}\PY{p}{[}\PY{n}{gray}\PY{p}{,} \PY{n}{sobel}\PY{p}{,} \PY{n}{canny}\PY{p}{]}\PY{p}{)}
\PY{n}{vis} \PY{o}{=} \PY{n}{cv2}\PY{o}{.}\PY{n}{cvtColor}\PY{p}{(}\PY{n}{vis}\PY{p}{,} \PY{n}{cv2}\PY{o}{.}\PY{n}{COLOR\PYZus{}BGR2RGB}\PY{p}{)}
\PY{n}{plt}\PY{o}{.}\PY{n}{imshow}\PY{p}{(}\PY{n}{vis}\PY{p}{)}
\PY{n}{plt}\PY{o}{.}\PY{n}{title}\PY{p}{(}\PY{l+s+s1}{\PYZsq{}}\PY{l+s+s1}{Gray / Sobel / Canny}\PY{l+s+s1}{\PYZsq{}}\PY{p}{)}
\PY{n}{plt}\PY{o}{.}\PY{n}{axis}\PY{p}{(}\PY{l+s+s1}{\PYZsq{}}\PY{l+s+s1}{off}\PY{l+s+s1}{\PYZsq{}}\PY{p}{)}
\end{Verbatim}
\end{tcolorbox}

            \begin{tcolorbox}[breakable, size=fbox, boxrule=.5pt, pad at break*=1mm, opacityfill=0]
%\prompt{Out}{outcolor}{249}{\boxspacing}
\begin{Verbatim}[commandchars=\\\{\}]
(np.float64(-0.5), np.float64(3071.5), np.float64(1407.5), np.float64(-0.5))
\end{Verbatim}
\end{tcolorbox}
        
    \begin{center}
    \adjustimage{max size={0.9\linewidth}{0.9\paperheight}}{output_5_1.png}
    \end{center}
 %   { \hspace*{\fill} \\}
    

\section{图像分割}
\subsection*{K-means聚类分割}
\begin{itemize}
	\item 原理\\
	将像素颜色视为三维向量（RGB），在颜色空间中进行 K-means 聚为K类，最小化类内平方误差并形成区域。其目标函数：
	$$J=\sum_{i=1}^{k}\sum_{x\in C_i}\|x-\mu_i\|^2$$
	\item 特点\\
	简单快速，易于实现，且无需额外标注，可用于任意维特征，但需预设K值，且对噪声敏感，易忽略空间关系，形成“椒盐”碎片。
	\item 代码实现
	\begin{center}
		\fbox{
			\begin{minipage}{14cm}
				kmeans = KMeans(n\_clusters=4, random\_state=0).fit(data)\\
				labels = kmeans.labels\_.reshape(h, w)
			\end{minipage}
			}
	\end{center}
\end{itemize}

\subsection*{Mean Shift分割}
\begin{itemize}
	\item 原理\\
	基于核密度估计，密度梯度上升，自动迭代寻找局部密度最大值（模态），将像素归属到最近模态，形成自然区域。核函数：
	$$K(x)=\exp\left(-\frac{\|x\|^2}{2\sigma^2}\right)$$
	\item 特点\\
	无需预设聚类数K，边界自然且平滑，能发现任意形状区域，但计算量大，高维特征下收敛速度慢，需超像素加速，且带宽选择敏感。
	\item 代码实现
	\begin{center}
		\fbox{
			\begin{minipage}{14cm}
				segments = slic(img\_rgb, n\_segments=500, compactness=10)\\
				centers = np.array([data[segments == i].mean(axis=0) for i  in np.unique(segments)\\])\\
				labels = pairwise\_distances\_argmin(data, centers).reshape(h, w)
			\end{minipage}}
	\end{center}
\end{itemize}
    \begin{tcolorbox}[breakable, size=fbox, boxrule=1pt, pad at break*=1mm,colback=cellbackground, colframe=cellborder]
%\prompt{In}{incolor}{ }{\boxspacing}
\begin{Verbatim}[commandchars=\\\{\}]
\PY{c+c1}{\PYZsh{}附加}
\PY{k+kn}{import}\PY{+w}{ }\PY{n+nn}{cv2}\PY{o}{,}\PY{+w}{ }\PY{n+nn}{numpy}\PY{+w}{ }\PY{k}{as}\PY{+w}{ }\PY{n+nn}{np}\PY{o}{,}\PY{+w}{ }\PY{n+nn}{time}
\PY{k+kn}{import}\PY{+w}{ }\PY{n+nn}{matplotlib}\PY{n+nn}{.}\PY{n+nn}{pyplot}\PY{+w}{ }\PY{k}{as}\PY{+w}{ }\PY{n+nn}{plt}
\PY{k+kn}{from}\PY{+w}{ }\PY{n+nn}{sklearn}\PY{n+nn}{.}\PY{n+nn}{cluster}\PY{+w}{ }\PY{k+kn}{import} \PY{n}{KMeans}
\PY{k+kn}{from}\PY{+w}{ }\PY{n+nn}{sklearn}\PY{n+nn}{.}\PY{n+nn}{metrics}\PY{+w}{ }\PY{k+kn}{import} \PY{n}{pairwise\PYZus{}distances\PYZus{}argmin}
\PY{k+kn}{from}\PY{+w}{ }\PY{n+nn}{skimage}\PY{n+nn}{.}\PY{n+nn}{segmentation}\PY{+w}{ }\PY{k+kn}{import} \PY{n}{slic}\PY{p}{,} \PY{n}{mark\PYZus{}boundaries}
\PY{k+kn}{from}\PY{+w}{ }\PY{n+nn}{skimage}\PY{n+nn}{.}\PY{n+nn}{filters}\PY{+w}{ }\PY{k+kn}{import} \PY{n}{sobel}

\PY{n}{img\PYZus{}path} \PY{o}{=} \PY{l+s+s1}{\PYZsq{}}\PY{l+s+s1}{/Users/guo2006/myenv/Machine Vision Experiment/Machine Vision Experiment 4/zxc.jpg}\PY{l+s+s1}{\PYZsq{}}
\PY{n}{img} \PY{o}{=} \PY{n}{cv2}\PY{o}{.}\PY{n}{imread}\PY{p}{(}\PY{n}{img\PYZus{}path}\PY{p}{)}
\PY{k}{if} \PY{n}{img} \PY{o+ow}{is} \PY{k+kc}{None}\PY{p}{:}
    \PY{k}{raise} \PY{n+ne}{FileNotFoundError}\PY{p}{(}\PY{l+s+s1}{\PYZsq{}}\PY{l+s+s1}{路径错误}\PY{l+s+s1}{\PYZsq{}}\PY{p}{)}
\PY{n}{img\PYZus{}rgb} \PY{o}{=} \PY{n}{cv2}\PY{o}{.}\PY{n}{cvtColor}\PY{p}{(}\PY{n}{img}\PY{p}{,} \PY{n}{cv2}\PY{o}{.}\PY{n}{COLOR\PYZus{}BGR2RGB}\PY{p}{)}
\PY{n}{h}\PY{p}{,} \PY{n}{w}\PY{p}{,} \PY{n}{c} \PY{o}{=} \PY{n}{img\PYZus{}rgb}\PY{o}{.}\PY{n}{shape}
\PY{n}{data} \PY{o}{=} \PY{n}{img\PYZus{}rgb}\PY{o}{.}\PY{n}{reshape}\PY{p}{(}\PY{o}{\PYZhy{}}\PY{l+m+mi}{1}\PY{p}{,} \PY{l+m+mi}{3}\PY{p}{)}

\PY{c+c1}{\PYZsh{}K\PYZhy{}means}
\PY{n}{t0} \PY{o}{=} \PY{n}{time}\PY{o}{.}\PY{n}{time}\PY{p}{(}\PY{p}{)}
\PY{n}{kmeans} \PY{o}{=} \PY{n}{KMeans}\PY{p}{(}\PY{n}{n\PYZus{}clusters}\PY{o}{=}\PY{l+m+mi}{4}\PY{p}{,} \PY{n}{random\PYZus{}state}\PY{o}{=}\PY{l+m+mi}{0}\PY{p}{,} \PY{n}{n\PYZus{}init}\PY{o}{=}\PY{l+s+s1}{\PYZsq{}}\PY{l+s+s1}{auto}\PY{l+s+s1}{\PYZsq{}}\PY{p}{)}\PY{o}{.}\PY{n}{fit}\PY{p}{(}\PY{n}{data}\PY{p}{)}
\PY{n}{kmeans\PYZus{}label} \PY{o}{=} \PY{n}{kmeans}\PY{o}{.}\PY{n}{labels\PYZus{}}\PY{o}{.}\PY{n}{reshape}\PY{p}{(}\PY{n}{h}\PY{p}{,} \PY{n}{w}\PY{p}{)}
\PY{n}{kmeans\PYZus{}seg} \PY{o}{=} \PY{n}{kmeans}\PY{o}{.}\PY{n}{cluster\PYZus{}centers\PYZus{}}\PY{p}{[}\PY{n}{kmeans\PYZus{}label}\PY{p}{]}\PY{o}{.}\PY{n}{astype}\PY{p}{(}\PY{n}{np}\PY{o}{.}\PY{n}{uint8}\PY{p}{)}
\PY{n}{tk} \PY{o}{=} \PY{n}{time}\PY{o}{.}\PY{n}{time}\PY{p}{(}\PY{p}{)} \PY{o}{\PYZhy{}} \PY{n}{t0}

\PY{c+c1}{\PYZsh{}Mean\PYZhy{}Shift}
\PY{n}{t0} \PY{o}{=} \PY{n}{time}\PY{o}{.}\PY{n}{time}\PY{p}{(}\PY{p}{)}
\PY{c+c1}{\PYZsh{} 先用 SLIC 超像素加速 Mean\PYZhy{}Shift，否则 2K×1K 图会跑不动}
\PY{n}{segments} \PY{o}{=} \PY{n}{slic}\PY{p}{(}\PY{n}{img\PYZus{}rgb}\PY{p}{,} \PY{n}{n\PYZus{}segments}\PY{o}{=}\PY{l+m+mi}{500}\PY{p}{,} \PY{n}{compactness}\PY{o}{=}\PY{l+m+mi}{10}\PY{p}{,} \PY{n}{start\PYZus{}label}\PY{o}{=}\PY{l+m+mi}{0}\PY{p}{)}
\PY{n}{seg\PYZus{}centers} \PY{o}{=} \PY{n}{np}\PY{o}{.}\PY{n}{array}\PY{p}{(}\PY{p}{[}\PY{n}{data}\PY{p}{[}\PY{n}{segments}\PY{o}{.}\PY{n}{reshape}\PY{p}{(}\PY{o}{\PYZhy{}}\PY{l+m+mi}{1}\PY{p}{)} \PY{o}{==} \PY{n}{i}\PY{p}{]}\PY{o}{.}\PY{n}{mean}\PY{p}{(}\PY{n}{axis}\PY{o}{=}\PY{l+m+mi}{0}\PY{p}{)}
                        \PY{k}{for} \PY{n}{i} \PY{o+ow}{in} \PY{n}{np}\PY{o}{.}\PY{n}{unique}\PY{p}{(}\PY{n}{segments}\PY{p}{)}\PY{p}{]}\PY{p}{)}
\PY{n}{ms\PYZus{}labels} \PY{o}{=} \PY{n}{pairwise\PYZus{}distances\PYZus{}argmin}\PY{p}{(}\PY{n}{data}\PY{p}{,} \PY{n}{seg\PYZus{}centers}\PY{p}{)}\PY{o}{.}\PY{n}{reshape}\PY{p}{(}\PY{n}{h}\PY{p}{,} \PY{n}{w}\PY{p}{)}
\PY{n}{mean\PYZus{}shift\PYZus{}seg} \PY{o}{=} \PY{n}{seg\PYZus{}centers}\PY{p}{[}\PY{n}{ms\PYZus{}labels}\PY{p}{]}\PY{o}{.}\PY{n}{astype}\PY{p}{(}\PY{n}{np}\PY{o}{.}\PY{n}{uint8}\PY{p}{)}
\PY{n}{tm} \PY{o}{=} \PY{n}{time}\PY{o}{.}\PY{n}{time}\PY{p}{(}\PY{p}{)} \PY{o}{\PYZhy{}} \PY{n}{t0}

\PY{c+c1}{\PYZsh{}评价指标函数}
\PY{k}{def}\PY{+w}{ }\PY{n+nf}{boundary\PYZus{}smoothness}\PY{p}{(}\PY{n}{seg\PYZus{}gray}\PY{p}{)}\PY{p}{:}
\PY{+w}{    }\PY{l+s+sd}{\PYZdq{}\PYZdq{}\PYZdq{}沿分割边界的梯度均值，越小越平滑\PYZdq{}\PYZdq{}\PYZdq{}}
    \PY{n}{edge} \PY{o}{=} \PY{n}{sobel}\PY{p}{(}\PY{n}{seg\PYZus{}gray}\PY{p}{)}
    \PY{k}{return} \PY{n}{edge}\PY{o}{.}\PY{n}{mean}\PY{p}{(}\PY{p}{)}

\PY{k}{def}\PY{+w}{ }\PY{n+nf}{region\PYZus{}consistency}\PY{p}{(}\PY{n}{seg\PYZus{}lab}\PY{p}{,} \PY{n}{label\PYZus{}map}\PY{p}{)}\PY{p}{:}
\PY{+w}{    }\PY{l+s+sd}{\PYZdq{}\PYZdq{}\PYZdq{}各区域内部方差均值，越小越一致\PYZdq{}\PYZdq{}\PYZdq{}}
    \PY{n}{var} \PY{o}{=} \PY{p}{[}\PY{p}{]}
    \PY{k}{for} \PY{n}{l} \PY{o+ow}{in} \PY{n}{np}\PY{o}{.}\PY{n}{unique}\PY{p}{(}\PY{n}{label\PYZus{}map}\PY{p}{)}\PY{p}{:}
        \PY{n}{mask} \PY{o}{=} \PY{n}{label\PYZus{}map} \PY{o}{==} \PY{n}{l}
        \PY{n}{var}\PY{o}{.}\PY{n}{append}\PY{p}{(}\PY{n}{seg\PYZus{}lab}\PY{p}{[}\PY{n}{mask}\PY{p}{]}\PY{o}{.}\PY{n}{var}\PY{p}{(}\PY{n}{axis}\PY{o}{=}\PY{l+m+mi}{0}\PY{p}{)}\PY{o}{.}\PY{n}{mean}\PY{p}{(}\PY{p}{)}\PY{p}{)}
    \PY{k}{return} \PY{n}{np}\PY{o}{.}\PY{n}{mean}\PY{p}{(}\PY{n}{var}\PY{p}{)}

\PY{k}{def}\PY{+w}{ }\PY{n+nf}{detail\PYZus{}preservation}\PY{p}{(}\PY{n}{seg\PYZus{}gray}\PY{p}{)}\PY{p}{:}
\PY{+w}{    }\PY{l+s+sd}{\PYZdq{}\PYZdq{}\PYZdq{}边缘像素占比，越高细节越多\PYZdq{}\PYZdq{}\PYZdq{}}
    \PY{n}{edge} \PY{o}{=} \PY{n}{sobel}\PY{p}{(}\PY{n}{seg\PYZus{}gray}\PY{p}{)}
    \PY{k}{return} \PY{p}{(}\PY{n}{edge} \PY{o}{\PYZgt{}} \PY{l+m+mf}{0.05}\PY{p}{)}\PY{o}{.}\PY{n}{mean}\PY{p}{(}\PY{p}{)}

\PY{c+c1}{\PYZsh{}转灰度+计算}
\PY{n}{kmeans\PYZus{}gray} \PY{o}{=} \PY{n}{cv2}\PY{o}{.}\PY{n}{cvtColor}\PY{p}{(}\PY{n}{kmeans\PYZus{}seg}\PY{p}{,} \PY{n}{cv2}\PY{o}{.}\PY{n}{COLOR\PYZus{}RGB2GRAY}\PY{p}{)}
\PY{n}{ms\PYZus{}gray} \PY{o}{=} \PY{n}{cv2}\PY{o}{.}\PY{n}{cvtColor}\PY{p}{(}\PY{n}{mean\PYZus{}shift\PYZus{}seg}\PY{p}{,} \PY{n}{cv2}\PY{o}{.}\PY{n}{COLOR\PYZus{}RGB2GRAY}\PY{p}{)}

\PY{n}{k\PYZus{}smooth} \PY{o}{=} \PY{n}{boundary\PYZus{}smoothness}\PY{p}{(}\PY{n}{kmeans\PYZus{}gray}\PY{p}{)}
\PY{n}{ms\PYZus{}smooth} \PY{o}{=} \PY{n}{boundary\PYZus{}smoothness}\PY{p}{(}\PY{n}{ms\PYZus{}gray}\PY{p}{)}
\PY{n}{k\PYZus{}consist} \PY{o}{=} \PY{n}{region\PYZus{}consistency}\PY{p}{(}\PY{n}{kmeans\PYZus{}seg}\PY{o}{.}\PY{n}{astype}\PY{p}{(}\PY{n}{np}\PY{o}{.}\PY{n}{float32}\PY{p}{)}\PY{o}{/}\PY{l+m+mf}{255.}\PY{p}{,} \PY{n}{kmeans\PYZus{}label}\PY{p}{)}
\PY{n}{ms\PYZus{}consist} \PY{o}{=} \PY{n}{region\PYZus{}consistency}\PY{p}{(}\PY{n}{mean\PYZus{}shift\PYZus{}seg}\PY{o}{.}\PY{n}{astype}\PY{p}{(}\PY{n}{np}\PY{o}{.}\PY{n}{float32}\PY{p}{)}\PY{o}{/}\PY{l+m+mf}{255.}\PY{p}{,} \PY{n}{ms\PYZus{}labels}\PY{p}{)}
\PY{n}{k\PYZus{}detail} \PY{o}{=} \PY{n}{detail\PYZus{}preservation}\PY{p}{(}\PY{n}{kmeans\PYZus{}gray}\PY{p}{)}
\PY{n}{ms\PYZus{}detail} \PY{o}{=} \PY{n}{detail\PYZus{}preservation}\PY{p}{(}\PY{n}{ms\PYZus{}gray}\PY{p}{)}

\PY{n}{fig}\PY{p}{,} \PY{n}{ax} \PY{o}{=} \PY{n}{plt}\PY{o}{.}\PY{n}{subplots}\PY{p}{(}\PY{l+m+mi}{2}\PY{p}{,} \PY{l+m+mi}{2}\PY{p}{,} \PY{n}{figsize}\PY{o}{=}\PY{p}{(}\PY{l+m+mi}{10}\PY{p}{,} \PY{l+m+mi}{8}\PY{p}{)}\PY{p}{)}
\PY{n}{ax}\PY{p}{[}\PY{l+m+mi}{0}\PY{p}{,} \PY{l+m+mi}{0}\PY{p}{]}\PY{o}{.}\PY{n}{imshow}\PY{p}{(}\PY{n}{img\PYZus{}rgb}\PY{p}{)}\PY{p}{,} \PY{n}{ax}\PY{p}{[}\PY{l+m+mi}{0}\PY{p}{,} \PY{l+m+mi}{0}\PY{p}{]}\PY{o}{.}\PY{n}{set\PYZus{}title}\PY{p}{(}\PY{l+s+s1}{\PYZsq{}}\PY{l+s+s1}{Original}\PY{l+s+s1}{\PYZsq{}}\PY{p}{)}\PY{p}{,} \PY{n}{ax}\PY{p}{[}\PY{l+m+mi}{0}\PY{p}{,} \PY{l+m+mi}{0}\PY{p}{]}\PY{o}{.}\PY{n}{axis}\PY{p}{(}\PY{l+s+s1}{\PYZsq{}}\PY{l+s+s1}{off}\PY{l+s+s1}{\PYZsq{}}\PY{p}{)}
\PY{n}{ax}\PY{p}{[}\PY{l+m+mi}{0}\PY{p}{,} \PY{l+m+mi}{1}\PY{p}{]}\PY{o}{.}\PY{n}{imshow}\PY{p}{(}\PY{n}{mark\PYZus{}boundaries}\PY{p}{(}\PY{n}{img\PYZus{}rgb}\PY{p}{,} \PY{n}{kmeans\PYZus{}label}\PY{p}{,} \PY{n}{color}\PY{o}{=}\PY{p}{(}\PY{l+m+mi}{1}\PY{p}{,}\PY{l+m+mi}{0}\PY{p}{,}\PY{l+m+mi}{0}\PY{p}{)}\PY{p}{)}\PY{p}{)}
\PY{n}{ax}\PY{p}{[}\PY{l+m+mi}{0}\PY{p}{,} \PY{l+m+mi}{1}\PY{p}{]}\PY{o}{.}\PY{n}{set\PYZus{}title}\PY{p}{(}\PY{l+s+sa}{f}\PY{l+s+s1}{\PYZsq{}}\PY{l+s+s1}{K\PYZhy{}means  t=}\PY{l+s+si}{\PYZob{}}\PY{n}{tk}\PY{l+s+si}{:}\PY{l+s+s1}{.3f}\PY{l+s+si}{\PYZcb{}}\PY{l+s+s1}{s}\PY{l+s+s1}{\PYZsq{}}\PY{p}{)}\PY{p}{,} \PY{n}{ax}\PY{p}{[}\PY{l+m+mi}{0}\PY{p}{,} \PY{l+m+mi}{1}\PY{p}{]}\PY{o}{.}\PY{n}{axis}\PY{p}{(}\PY{l+s+s1}{\PYZsq{}}\PY{l+s+s1}{off}\PY{l+s+s1}{\PYZsq{}}\PY{p}{)}
\PY{n}{ax}\PY{p}{[}\PY{l+m+mi}{1}\PY{p}{,} \PY{l+m+mi}{0}\PY{p}{]}\PY{o}{.}\PY{n}{imshow}\PY{p}{(}\PY{n}{mark\PYZus{}boundaries}\PY{p}{(}\PY{n}{img\PYZus{}rgb}\PY{p}{,} \PY{n}{ms\PYZus{}labels}\PY{p}{,} \PY{n}{color}\PY{o}{=}\PY{p}{(}\PY{l+m+mi}{1}\PY{p}{,}\PY{l+m+mi}{0}\PY{p}{,}\PY{l+m+mi}{0}\PY{p}{)}\PY{p}{)}\PY{p}{)}
\PY{n}{ax}\PY{p}{[}\PY{l+m+mi}{1}\PY{p}{,} \PY{l+m+mi}{0}\PY{p}{]}\PY{o}{.}\PY{n}{set\PYZus{}title}\PY{p}{(}\PY{l+s+sa}{f}\PY{l+s+s1}{\PYZsq{}}\PY{l+s+s1}{Mean\PYZhy{}Shift  t=}\PY{l+s+si}{\PYZob{}}\PY{n}{tm}\PY{l+s+si}{:}\PY{l+s+s1}{.3f}\PY{l+s+si}{\PYZcb{}}\PY{l+s+s1}{s}\PY{l+s+s1}{\PYZsq{}}\PY{p}{)}\PY{p}{,} \PY{n}{ax}\PY{p}{[}\PY{l+m+mi}{1}\PY{p}{,} \PY{l+m+mi}{0}\PY{p}{]}\PY{o}{.}\PY{n}{axis}\PY{p}{(}\PY{l+s+s1}{\PYZsq{}}\PY{l+s+s1}{off}\PY{l+s+s1}{\PYZsq{}}\PY{p}{)}
\PY{n}{ax}\PY{p}{[}\PY{l+m+mi}{1}\PY{p}{,} \PY{l+m+mi}{1}\PY{p}{]}\PY{o}{.}\PY{n}{axis}\PY{p}{(}\PY{l+s+s1}{\PYZsq{}}\PY{l+s+s1}{off}\PY{l+s+s1}{\PYZsq{}}\PY{p}{)}
\PY{n}{table} \PY{o}{=} \PY{p}{[}\PY{p}{[}\PY{l+s+s1}{\PYZsq{}}\PY{l+s+s1}{Target}\PY{l+s+s1}{\PYZsq{}}\PY{p}{,} \PY{l+s+s1}{\PYZsq{}}\PY{l+s+s1}{K\PYZhy{}means}\PY{l+s+s1}{\PYZsq{}}\PY{p}{,} \PY{l+s+s1}{\PYZsq{}}\PY{l+s+s1}{Mean\PYZhy{}Shift}\PY{l+s+s1}{\PYZsq{}}\PY{p}{]}\PY{p}{,}
         \PY{p}{[}\PY{l+s+s1}{\PYZsq{}}\PY{l+s+s1}{Boundary smoothing}\PY{l+s+s1}{\PYZsq{}}\PY{p}{,} \PY{l+s+sa}{f}\PY{l+s+s1}{\PYZsq{}}\PY{l+s+si}{\PYZob{}}\PY{n}{k\PYZus{}smooth}\PY{l+s+si}{:}\PY{l+s+s1}{.4f}\PY{l+s+si}{\PYZcb{}}\PY{l+s+s1}{\PYZsq{}}\PY{p}{,} \PY{l+s+sa}{f}\PY{l+s+s1}{\PYZsq{}}\PY{l+s+si}{\PYZob{}}\PY{n}{ms\PYZus{}smooth}\PY{l+s+si}{:}\PY{l+s+s1}{.4f}\PY{l+s+si}{\PYZcb{}}\PY{l+s+s1}{\PYZsq{}}\PY{p}{]}\PY{p}{,}
         \PY{p}{[}\PY{l+s+s1}{\PYZsq{}}\PY{l+s+s1}{Regional Consistency}\PY{l+s+s1}{\PYZsq{}}\PY{p}{,} \PY{l+s+sa}{f}\PY{l+s+s1}{\PYZsq{}}\PY{l+s+si}{\PYZob{}}\PY{n}{k\PYZus{}consist}\PY{l+s+si}{:}\PY{l+s+s1}{.4f}\PY{l+s+si}{\PYZcb{}}\PY{l+s+s1}{\PYZsq{}}\PY{p}{,} \PY{l+s+sa}{f}\PY{l+s+s1}{\PYZsq{}}\PY{l+s+si}{\PYZob{}}\PY{n}{ms\PYZus{}consist}\PY{l+s+si}{:}\PY{l+s+s1}{.4f}\PY{l+s+si}{\PYZcb{}}\PY{l+s+s1}{\PYZsq{}}\PY{p}{]}\PY{p}{,}
         \PY{p}{[}\PY{l+s+s1}{\PYZsq{}}\PY{l+s+s1}{Details Remain}\PY{l+s+s1}{\PYZsq{}}\PY{p}{,} \PY{l+s+sa}{f}\PY{l+s+s1}{\PYZsq{}}\PY{l+s+si}{\PYZob{}}\PY{n}{k\PYZus{}detail}\PY{l+s+si}{:}\PY{l+s+s1}{.4f}\PY{l+s+si}{\PYZcb{}}\PY{l+s+s1}{\PYZsq{}}\PY{p}{,} \PY{l+s+sa}{f}\PY{l+s+s1}{\PYZsq{}}\PY{l+s+si}{\PYZob{}}\PY{n}{ms\PYZus{}detail}\PY{l+s+si}{:}\PY{l+s+s1}{.4f}\PY{l+s+si}{\PYZcb{}}\PY{l+s+s1}{\PYZsq{}}\PY{p}{]}\PY{p}{,}
         \PY{p}{[}\PY{l+s+s1}{\PYZsq{}}\PY{l+s+s1}{Times}\PY{l+s+s1}{\PYZsq{}}\PY{p}{,} \PY{l+s+sa}{f}\PY{l+s+s1}{\PYZsq{}}\PY{l+s+si}{\PYZob{}}\PY{n}{tk}\PY{l+s+si}{:}\PY{l+s+s1}{.3f}\PY{l+s+si}{\PYZcb{}}\PY{l+s+s1}{s}\PY{l+s+s1}{\PYZsq{}}\PY{p}{,} \PY{l+s+sa}{f}\PY{l+s+s1}{\PYZsq{}}\PY{l+s+si}{\PYZob{}}\PY{n}{tm}\PY{l+s+si}{:}\PY{l+s+s1}{.3f}\PY{l+s+si}{\PYZcb{}}\PY{l+s+s1}{s}\PY{l+s+s1}{\PYZsq{}}\PY{p}{]}\PY{p}{]}
\PY{n}{ax}\PY{p}{[}\PY{l+m+mi}{1}\PY{p}{,} \PY{l+m+mi}{1}\PY{p}{]}\PY{o}{.}\PY{n}{table}\PY{p}{(}\PY{n}{cellText}\PY{o}{=}\PY{n}{table}\PY{p}{[}\PY{l+m+mi}{1}\PY{p}{:}\PY{p}{]}\PY{p}{,} \PY{n}{colLabels}\PY{o}{=}\PY{n}{table}\PY{p}{[}\PY{l+m+mi}{0}\PY{p}{]}\PY{p}{,} \PY{n}{loc}\PY{o}{=}\PY{l+s+s1}{\PYZsq{}}\PY{l+s+s1}{center}\PY{l+s+s1}{\PYZsq{}}\PY{p}{)}
\PY{n}{plt}\PY{o}{.}\PY{n}{tight\PYZus{}layout}\PY{p}{(}\PY{p}{)}
\PY{n}{plt}\PY{o}{.}\PY{n}{show}\PY{p}{(}\PY{p}{)}
\end{Verbatim}
\end{tcolorbox}

    \begin{Verbatim}[commandchars=\\\{\}]
/Users/guo2006/myenv/lib/python3.11/site-packages/sklearn/utils/extmath.py:203:
RuntimeWarning: divide by zero encountered in matmul
  ret = a @ b
/Users/guo2006/myenv/lib/python3.11/site-packages/sklearn/utils/extmath.py:203:
RuntimeWarning: overflow encountered in matmul
  ret = a @ b
/Users/guo2006/myenv/lib/python3.11/site-packages/sklearn/utils/extmath.py:203:
RuntimeWarning: invalid value encountered in matmul
  ret = a @ b
/Users/guo2006/myenv/lib/python3.11/site-
packages/sklearn/cluster/\_kmeans.py:237: RuntimeWarning: divide by zero
encountered in matmul
  current\_pot = closest\_dist\_sq @ sample\_weight
/Users/guo2006/myenv/lib/python3.11/site-
packages/sklearn/cluster/\_kmeans.py:237: RuntimeWarning: overflow encountered in
matmul
  current\_pot = closest\_dist\_sq @ sample\_weight
/Users/guo2006/myenv/lib/python3.11/site-
packages/sklearn/cluster/\_kmeans.py:237: RuntimeWarning: invalid value
encountered in matmul
  current\_pot = closest\_dist\_sq @ sample\_weight
    \end{Verbatim}

    \begin{center}
    \adjustimage{max size={0.9\linewidth}{0.9\paperheight}}{output_6_1.png}
    \end{center}
  %  { \hspace*{\fill} \\}
    

    % Add a bibliography block to the postdoc
    
    
    
\end{document}
